\documentclass[oneside,a4papar,11pt,UTF8]{book}
%————————————————————————————————————————————————————————————————————————————————————————
%%%%%%%%%%%%%%%%%%%%%%%%%%%%这里调用了曲豆豆的独家秘籍%%%%%%%%%%%%%%%%%%%%%%%%%%%%%%%%%%%%%%%%
\usepackage{ctex}
\usepackage{amsfonts,amsmath,amscd,amssymb,amsthm} %美国数学会套装
\usepackage{latexsym,amsbsy,bbold, fullpage}
\usepackage{pdfsync,yfonts}
\usepackage[all,pdf]{xy}
\usepackage{tkz-graph,tikz-cd,xy}
\usepackage{color,xcolor,tcolorbox,framed}      %颜色
\usepackage{fancyhdr,lastpage}                  %页眉页脚，最后一页
\usepackage{incgraph}
\usepackage{multicol}                           %双栏
\usepackage{enumerate}                          %枚举

\usepackage{geometry}
\newgeometry{left = 3.5 cm, right=3.5cm}

\usepackage[T1]{fontenc}                       %Required for accented characters
\usepackage{titlesec}                          %章节标题风格
\usepackage{hyperref,url}                      %超链接
\usepackage{mathpazo}                          %曲豆豆喜欢用的Palatino字体
\usepackage[framemethod=tikz]{mdframed}        %花里胡哨的框框
\usepackage{wallpaper}                         %背景图片
\usepackage{pstricks,pst-node,pst-tree}        %图论（暂时没卵用）
\usepackage{graphicx,wrapfig,float}            % 用来插图
\usepackage{ifthen}                            %逻辑判断
\usepackage{tikz,stmaryrd}                     %究极绘图语言
\usepackage{booktabs}                          %调整表格线与上下内容的间隔
\usepackage{chngpage}
\usepackage{enumitem}
\usepackage{enumerate}
\usepackage{makeidx}
\usepackage[tikz]{bclogo}

\usepackage{microtype}
%\usepackage[american]{babel}

\linespread{1.2} % 调整行间距
\allowdisplaybreaks[4]%允许跨页公式
%————————————————————————————————————————————————————

\newenvironment{solution}
  {\begin{proof}[解]\,}
  {\end{proof}}

%————————————————————————————————————————————————————

\newenvironment{prob}
  {\item\bfseries\kaishu}
  {}

\renewcommand{\proofname}{证明}

%-------------------------------------------------------------

\newcommand*{\Gray}[1]{\textcolor[rgb]{0.62,0.59,0.56}{#1}}

%——————————————————————————————————————————————————————————————————

\everymath{\displaystyle}  %全部行间公式，暗黑魔咒，慎用！

\renewcommand{\labelenumi}{\textbf{\theenumi}.\,}  %枚举环境编号设置
%\setlist[enumerate,2]{label=\arabic}
\renewcommand{\labelenumii}{({\theenumii})\,}  %枚举环境编号设置

%——————————索引标签术——————————————————%

\newcommand*{\newdef}[2]{
    \textbf{#1}(#2)
    \index{#2 $\quad$ #1}
    \label{def-#1}
  }

%——————————Hint———————————————————%

\newcommand*{\hint}[1]
  {
{\fangsong\textcolor[rgb]{0.6,0.6,0.6}{[提示:#1]}}
  }

\newcommand*{\remark}[1]
  {
    ({\fangsong #1})
  } 
                      %排版命令，重要！
%这是曲豆豆的红包
\usepackage{extarrows}   %更多箭头

\newcommand*{\vs}{\vspace{5pt}}
\newcommand*{\vsp}{\vspace{10pt}}
\newcommand*{\vspp}{\vspace{20pt}}
\newcommand*{\kong}{$\,\,\,$}

%%%%%%%%%%%%%%%%%%%%%%%%%%%%%%%%%%%%%%%%%%%%%%%%%%%%%%%%%%%%
\newcommand*{\Langle}{\left\langle}
\newcommand*{\Rangle}{\right\rangle}
%%%%%%%%%%%%%%%%%%%%%%%%%%%%%%%%%%%%%%%%%%%%%%%%%%%%%%%%%%%%
\newcommand*{\bbA}{\mathbb{A}}
\newcommand*{\bbB}{\mathbb{B}}
\newcommand*{\bbC}{\mathbb{C}}
\newcommand*{\bbD}{\mathbb{D}}
\newcommand*{\bbE}{\mathbb{E}}
\newcommand*{\bbF}{\mathbb{F}}
\newcommand*{\bbG}{\mathbb{G}}
\newcommand*{\bbH}{\mathbb{H}}
\newcommand*{\bbI}{\mathbb{I}}
\newcommand*{\bbJ}{\mathbb{J}}
\newcommand*{\bbK}{\mathbb{K}}
\newcommand*{\bbL}{\mathbb{L}}
\newcommand*{\bbM}{\mathbb{M}}
\newcommand*{\bbN}{\mathbb{N}}
\newcommand*{\bbO}{\mathbb{O}}
\newcommand*{\bbP}{\mathbb{P}}
\newcommand*{\bbQ}{\mathbb{Q}}
\newcommand*{\bbR}{\mathbb{R}}
\newcommand*{\bbS}{\mathbb{S}}
\newcommand*{\bbT}{\mathbb{T}}
\newcommand*{\bbU}{\mathbb{U}}
\newcommand*{\bbV}{\mathbb{V}}
\newcommand*{\bbW}{\mathbb{W}}
\newcommand*{\bbX}{\mathbb{X}}
\newcommand*{\bbY}{\mathbb{Y}}
\newcommand*{\bbZ}{\mathbb{Z}}

\newcommand*{\mcalA}{\mathcal{A}}
\newcommand*{\mcalB}{\mathcal{B}}
\newcommand*{\mcalC}{\mathcal{C}}
\newcommand*{\mcalD}{\mathcal{D}}
\newcommand*{\mcalE}{\mathcal{E}}
\newcommand*{\mcalF}{\mathcal{F}}
\newcommand*{\mcalG}{\mathcal{G}}
\newcommand*{\mcalH}{\mathcal{H}}
\newcommand*{\mcalI}{\mathcal{I}}
\newcommand*{\mcalJ}{\mathcal{J}}
\newcommand*{\mcalK}{\mathcal{K}}
\newcommand*{\mcalL}{\mathcal{L}}
\newcommand*{\mcalM}{\mathcal{M}}
\newcommand*{\mcalN}{\mathcal{N}}
\newcommand*{\mcalO}{\mathcal{O}}
\newcommand*{\mcalP}{\mathcal{P}}
\newcommand*{\mcalQ}{\mathcal{Q}}
\newcommand*{\mcalR}{\mathcal{R}}
\newcommand*{\mcalS}{\mathcal{S}}
\newcommand*{\mcalT}{\mathcal{T}}
\newcommand*{\mcalU}{\mathcal{U}}
\newcommand*{\mcalV}{\mathcal{V}}
\newcommand*{\mcalW}{\mathcal{W}}
\newcommand*{\mcalX}{\mathcal{X}}
\newcommand*{\mcalY}{\mathcal{Y}}
\newcommand*{\mcalZ}{\mathcal{Z}}

\newcommand*{\mfka}{\mathfrak{a}}
\newcommand*{\mfkg}{\mathfrak{g}}    %Lie algebra
\newcommand*{\mfkh}{\mathfrak{h}}
\newcommand*{\mfkm}{\mathfrak{m}}    %maximal ideal
\newcommand*{\mfkp}{\mathfrak{p}}    %prime ideal
\newcommand*{\mfkq}{\mathfrak{q}}

\newcommand*{\mfkgl}{\mathfrak{gl}}
\newcommand*{\mfksl}{\mathfrak{sl}}
\newcommand*{\mfkso}{\mathfrak{so}}
\newcommand*{\mfksp}{\mathfrak{sp}}

\newcommand*{\bfa}{\boldsymbol{a}}
\newcommand*{\bfe}{\boldsymbol{e}}
\newcommand*{\bfk}{\boldsymbol{k}}
\newcommand*{\bfl}{\boldsymbol{l}}
\newcommand*{\bfm}{\boldsymbol{m}}
\newcommand*{\bfn}{\boldsymbol{n}}
\newcommand*{\bfp}{\boldsymbol{p}}
\newcommand*{\bfq}{\boldsymbol{q}}
\newcommand*{\bfr}{\boldsymbol{r}}
\newcommand*{\bfu}{\boldsymbol{u}}
\newcommand*{\bfv}{\boldsymbol{v}}
\newcommand*{\bfw}{\boldsymbol{w}}
\newcommand*{\bfx}{\boldsymbol{x}}
\newcommand*{\bfy}{\boldsymbol{y}}
\newcommand*{\bfz}{\boldsymbol{z}}

\newcommand*{\bfF}{\boldsymbol{F}}
\newcommand*{\bfS}{\boldsymbol{S}}
\newcommand*{\bfrho}{\boldsymbol{\rho}}
\newcommand*{\bfveps}{\boldsymbol{\varepsilon}}

\newcommand*{\uhat}{\widehat{u}}
\newcommand*{\vhat}{\widehat{v}}

\newcommand*{\Ahat}{\widehat{A}}
\newcommand*{\Bhat}{\widehat{B}}
\newcommand*{\Chat}{\widehat{C}}
\newcommand*{\Dhat}{\widehat{D}}
\newcommand*{\Ehat}{\widehat{E}}
\newcommand*{\Fhat}{\widehat{F}}
\newcommand*{\Ghat}{\widehat{G}}
\newcommand*{\Hhat}{\widehat{H}}

\newcommand*{\zbar}{\overline{z}}            %complex number
%\newcommand*{\wbar}{\overline{w}}
\newcommand*{\taubar}{\overline{\tau}}

\newcommand*{\ftil}{\widetilde{f}}
\newcommand*{\gtil}{\widetilde{g}}
\newcommand*{\htil}{\widetilde{h}}
\newcommand*{\itil}{\widetilde{i}}
\newcommand*{\xtil}{\widetilde{x}}
\newcommand*{\ytil}{\widetilde{y}}
\newcommand*{\ztil}{\widetilde{z}}
\newcommand*{\wtil}{\widetilde{w}}

\newcommand*{\Ftil}{\widetilde{F}}
\newcommand*{\Gtil}{\widetilde{G}}
\newcommand*{\Htil}{\widetilde{H}}
\newcommand*{\Ntil}{\widetilde{N}}

\newcommand*{\faitil}{\widetilde{\varphi}}
\newcommand*{\psitil}{\widetilde{\psi}}
\newcommand*{\pitil}{\widetilde{\pi}}
\newcommand*{\Omgtil}{\widetilde{\Omega}}

\newcommand*{\ola}{\overline{a}}
\newcommand*{\olb}{\overline{b}}
\newcommand*{\olc}{\overline{c}}
\newcommand*{\old}{\overline{d}}
\newcommand*{\ole}{\overline{e}}
\newcommand*{\olf}{\overline{f}}
\newcommand*{\olg}{\overline{g}}
\newcommand*{\olh}{\overline{h}}
\newcommand*{\oli}{\overline{i}}
\newcommand*{\olj}{\overline{j}}
\newcommand*{\olk}{\overline{k}}
\newcommand*{\oll}{\overline{l}}
\newcommand*{\olm}{\overline{m}}
\newcommand*{\oln}{\overline{n}}
\newcommand*{\olo}{\overline{o}}
\newcommand*{\olp}{\overline{p}}
\newcommand*{\olq}{\overline{q}}
\newcommand*{\olr}{\overline{r}}
\newcommand*{\ols}{\overline{s}}
\newcommand*{\olt}{\overline{t}}
\newcommand*{\olu}{\overline{u}}
\newcommand*{\olv}{\overline{v}}
\newcommand*{\olw}{\overline{w}}
\newcommand*{\olx}{\overline{x}}
\newcommand*{\oly}{\overline{y}}
\newcommand*{\olz}{\overline{z}}

\newcommand*{\olA}{\overline{A}}
\newcommand*{\olB}{\overline{B}}
\newcommand*{\olG}{\overline{G}}

\newcommand*{\afa}{\alpha}
\newcommand*{\lmd}{\lambda}
\newcommand*{\Lmd}{\Lambda}
\newcommand*{\sgm}{\sigma}
\newcommand*{\Sgm}{\Sigma}
\newcommand*{\gma}{\gamma}
\newcommand*{\Gma}{\Gamma}
\newcommand*{\dta}{\delta}
\newcommand*{\Dta}{\Delta}
\newcommand*{\omg}{\omega}
\newcommand*{\Omg}{\Omega}
\newcommand*{\veps}{\varepsilon}
\newcommand*{\fai}{\varphi}
\newcommand*{\pai}{\varpi}

%%%%%%%%%%%%%%%%%%%%%%%%%%%%%%%%%%%%%%%%%%%%%%%%

\newcommand*{\ra}{\rightarrow}
\newcommand*{\inj}{\hookrightarrow}
\newcommand*{\surj}{\twoheadrightarrow}
\newcommand*{\xra}{\xrightarrow}
\newcommand*{\xla}{\xleftarrow}
\newcommand*{\xle}{\xlongequal}
\newcommand*{\llra}{\longleftrightarrow}
\newcommand*{\Llra}{\Longleftrightarrow}
\newcommand*{\rlim}{\varinjlim}                 %colimit，余极限，归纳极限，正向极限
\newcommand*{\llim}{\varprojlim}                %limit，极限，投射极限，逆向极限
\newcommand*{\suobing}{\lrcorner\,}             %张量的缩并
\newcommand*{\qqquad}{\quad\,\quad\,\quad}

\newcommand*{\xybigrow}{\xymatrixrowsep{5pc}}   %交换图排版专用
\newcommand*{\xybigcol}{\xymatrixcolsep{5pc}}   %增大行、列间距
\newcommand*{\tabularbigrow}{\specialrule{0em}{4pt}{4pt}}

\newcommand*{\pz}{\text{---}}                   %破折号
\newcommand*{\yc}{\triangle}                    %余代数的乘法，co-product
\newcommand*{\p}{\partial}                      %偏微分、边缘算子partial
\newcommand*{\td}{\,\mathrm{d}}                   %微分算子d
\newcommand*{\pbar}{\overline{\partial}}        %Dolbeaut算子partial-bar
\newcommand*{\ten}{\otimes}                     %张量积tensor product（改为算子形式）
\newcommand*{\bigten}{\bigotimes}
\newcommand*{\jing}{\#\,}
\newcommand*{\op}{^{\text{op}}}
\newcommand*{\downdot}{_{\bullet}}              %下方黑点，用于同调代数中的链复形
\newcommand*{\ddowndot}{_{\bullet\bullet}}
\newcommand*{\updot}{^{\bullet}}
\newcommand*{\wedgeform}{\bigwedge\nolimits^}
\newcommand*{\ssubset}{\subset\!\subset}        %紧包含
\newcommand*{\vkong}{\varnothing}               %空集

%%%%%%%%%%%%%%%%%%%%%%%%%%%%%%%%%%%%%%%%%%%%%%%%

\newcommand*{\lbar}[1]{\overline{#1}}    %比较长的上横线

\newcommand*{\fps}[1]{[\![#1]\!]}        %Formal power series形式幂级数
\newcommand*{\Ls}[1]{(\!(#1)\!)}         %Laurent series     洛朗级数

\newcommand*{\Bigset}[2]
  {
   \Big\{ #1 \Big| #2 \Big\}             %大集合符号
  }

\newcommand*{\bigset}[2]
  {
   \big\{ #1 \big| #2 \big\}             %较大的集合符号
  }

\newcommand*{\pair}[2]                   %普通的内积
  {
    \langle
      #1,#2
    \rangle
  }

\newcommand*{\Pair}[2]                   %大的的内积
  {
    \left\langle
      #1,#2
    \right\rangle
  }

\newcommand*{\ppair}[2]                   %双括号内积
  {
    \langle\!\!
      \langle
        #1,#2
      \rangle\!\!
    \rangle
  }

\newcommand*{\Ppair}[2]                   %大双括号内积
  {
    \left\langle\!\!
      \left\langle
        #1,#2
      \right\rangle\!\!
    \right\rangle
  }

\newcommand*{\norm}[1]                   %范数
  {
    |\!|#1|\!|
  }

\newcommand*{\Norm}[1]                   %大型范数
  {
    \left|\!
      \left|
        #1
      \right|\!
    \right|
  }

\newcommand*{\nnorm}[1]                   %三竖线范数
  {
    |\!|\!|
      #1
    |\!|\!|
  }

\newcommand*{\Nnorm}[1]                   %大型三竖线范数
  {
    \left|\!
      \left|\!
        \left|
          #1
        \right|\!
      \right|\!
    \right|
  }

\newcommand*{\Choose}[2]                  %方括号组合数
  {\left[{
      #1 \atop
      #2
  }\right]}

\newcommand*{\dd}[1]                       %常微分算子
  {\frac{\mathrm{d}   }
        {\mathrm{d} #1}
  }

\newcommand*{\ddfrac}[2]
  {\frac{\mathrm{d} #1}
        {\mathrm{d} #2}
  }                                       %一阶常微分

\newcommand*{\pp}[1]
  {\frac{\partial   }
        {\partial #1}
  }                                       %切向量场，一阶偏微分算子

\newcommand*{\pfrac}[2]
  {\frac{\partial #1}
        {\partial #2}
  }                                       %一阶偏微分

\newcommand*{\ppfrac}[2]
  {\frac{\partial^2{#1}}
        {\partial{#2}^2}}                 %二阶偏微分

\newcommand*{\pmfrac}[3]
  {\frac{\partial^2{#1}}
        {\partial{#2}\partial{#3}}}        %二阶混合偏微分

\newcommand*{\mulint}[1]                   %多重积分
  {
    \underbrace{
      \int\int\cdots\int
    }_{#1}\,\!
  }

%amsmath 宏包输入矩阵%
%matrix pmatrix bmatrix Bmatrix vmatrix Vmatrix%
%分别为无括号，小括号，中括号，大括号，竖线，双竖线%
%%%%%%%%%%%%%%%%%%%%%%%%%%%%%%%%%%%%%%%%%%%%%%%%
\DeclareMathOperator{\ad}{ad}         %adjoint
\DeclareMathOperator{\Ann}{Ann}         %adjoint
\DeclareMathOperator{\Ass}{Ass}       %Associated Prime
\DeclareMathOperator{\Aut}{Aut}       %automorphism
\DeclareMathOperator{\CE}{CE}         %Chevalley-Eilenberg
\DeclareMathOperator{\ch}{char}
\DeclareMathOperator{\cl}{cl}         %classical
\DeclareMathOperator{\Coder}{Coder}   %Co-derivation  并不是“码农”的意思
\DeclareMathOperator{\coker}{coker}   %co-kernel      余核
\DeclareMathOperator{\Conf}{Conf}     %Kontsevich公式当中用到了；configuration of points
\DeclareMathOperator{\Crit}{Crit}     %Critical locus
\DeclareMathOperator{\Der}{Der}       %Derivation        导子
\DeclareMathOperator{\Div}{div}       %divergence     散度
\DeclareMathOperator{\DR}{DR}         %De-Rham
\DeclareMathOperator{\End}{End}       %Endomorphism
\DeclareMathOperator{\Ext}{Ext}
\DeclareMathOperator{\ev}{ev}         %evaluation
\DeclareMathOperator{\Free}{Free}
\DeclareMathOperator{\Gal}{Gal}       %Galois
\DeclareMathOperator{\GL}{GL}         %General linear
\DeclareMathOperator{\HC}{HC}         %Hochschild Cyclic (co)homology
\DeclareMathOperator{\Hess}{Hess}     %Hessian
\DeclareMathOperator{\HH}{HH}         %Hochschild (co)homology
\DeclareMathOperator{\Hom}{Hom}       %Homorphism
\DeclareMathOperator{\id}{id}         %identity
\DeclareMathOperator{\im}{Im}         %Image
\DeclareMathOperator{\Inn}{Inn}       %Inner Derivation  内导子
\DeclareMathOperator{\Inv}{Inv}       %Invariant
\DeclareMathOperator{\lcm}{l.c.m.}    %least commom ...最小公倍数
\DeclareMathOperator{\Lie}{Lie}       %Lie algebra or Lie derivation
\DeclareMathOperator{\LT}{LT}         %Leading Term
\DeclareMathOperator{\Map}{Map}       %Mapping
\DeclareMathOperator{\Mor}{Mor}       %Morphism
\DeclareMathOperator{\nil}{nil}       %nilradical
\DeclareMathOperator{\Obj}{Obj}       %Object
\DeclareMathOperator{\Obs}{Obs}       %Observable
\DeclareMathOperator{\per}{per}       %periodic cyclic complex
\DeclareMathOperator{\PV}{PV}         %Polyvector field
\DeclareMathOperator{\rad}{rad}       %radical
\DeclareMathOperator{\ran}{ran}       %range
\DeclareMathOperator{\rank}{rank}     %rank 秩
\DeclareMathOperator{\Rep}{Rep}
\DeclareMathOperator{\sgn}{sgn}
\DeclareMathOperator{\Sh}{Sh}
\DeclareMathOperator{\SL}{SL}         %special linear
\DeclareMathOperator{\Span}{span}
\DeclareMathOperator{\Spec}{Spec}     %素谱
\DeclareMathOperator{\supp}{supp}     %Support 支撑集
\DeclareMathOperator{\Sym}{Sym}       %Symmetry  对称张量积
\DeclareMathOperator{\Tor}{Tor}
\DeclareMathOperator{\Tot}{Tot}       %Total complex of a double-complex
\DeclareMathOperator{\trdeg}{trdeg}   %transcendence degree
\DeclareMathOperator{\YM}{YM}         %Yang-Mills   杨振宁-米尔斯

%%%%%%%%%%%%%%%%%%%%%%%%%%%%%%%%%%%%%%%%%%%%%%%%
                       %数学公式输入法
%%%%%%%%%%%%%%%%%%%%%%%%%%%%% 禁 忌 之 地，非 请 勿 进 %%%%%%%%%%%%%%%%%%%%%%%%%%%%%%%%%%%%%%%%%%
%————————————————————————————————————————————————————————————————————————————————————————
\begin{document}

\tableofcontents
\newpage

\chapter{模板使用说明}

\section{总则}

本模板供\textbf{曲豆豆的高等数学群}
成员自愿帮曲豆豆编写\textbf{代数学千题解}使用.
目的在于提升志愿者的代数学功底,
兼提高\LaTeX 操作水平.

\begin{enumerate}
  \item Dummit书中的每一节分为
  \textbf{内容概要}与\textbf{习题与解答}两部分.
  
  \item \textbf{内容概要}部分
  为Dummit书中该节内容的总结.
  撰写此部分要注意:
  \begin{enumerate}
    \item 撰写这一部分之前，
    请先将Dummit书的原文读透,
    保证理解理解、掌握书中的所有知识点.
    
    \item 撰写这一部分时，
    要将这节出现的概念、记号交代清楚，
    将重要定理、重要例子简要介绍
    \remark{不必整理定理证明;
    不必所有的小例子,小推论面面俱到,
    请根据自己的理解与判断,提炼出主要的内容.}
     
    \item 不要逐字翻译原书,
    而要凭自己的理解,用自己的话来表述.
    
    \item 专业词汇的中文翻译,
    请与曲豆豆沟通.不建议自作主张. 
    
    \item 排版规范: 请使用
    \verb"itemize" 枚举环境
    将各个要点逐条列举.
    详见示例文本.
    
    \item 若该节内容较多,可用
    \verb"subsubsection"来划分更细的段落层次,
    每个\verb"subsubsection"的标题自拟.
  \end{enumerate}
  
  \item \textbf{习题与解答}部分
  是代数学千题解的主体,
  内容包括该节的全部习题的解答.
  撰写此部分要注意:
  \begin{enumerate}
    \item 全部习题,解答必须位于枚举环境
    \verb"enumerate"当中.
    
    \item 在\verb"prob"环境中输入习题,
    该环境里可以再嵌套枚举环境
    \verb"enumerate"来写此题目的各个小问.
    
    \item 提供证明环境\verb"proof"
    与解答环境\verb"solution".
    以上环境的具体用法详见示例文本.
    
    \item 提供含参命令\verb"hint",
    用于题目后面的\textbf{提示}的排版.
    
    用法为\verb"\hint{提示的内容写到这里}".
    
    \item 提供含参命令\verb"remark",
    用于题目的\textbf{注记}排版,
    也可用于吐槽.

    用法为\verb"\remark{吐槽的内容写到这里}".
  \end{enumerate}
\end{enumerate}

\section{文字书写规范}

假定大家了解\LaTeX 公式输入的基础知识.
本节介绍代数学千题解模板当中的特殊注意事项:


\begin{enumerate}
  \item 禁止使用中文标点符号!全部使用西文标点.
  
  \item 关于数学字体.
  在代数学千题解当中,
  默认的数学公式字体为\verb"mathpazo".
  此外,罗马体,空心体,手写体,哥特体,无衬线体,粗体
  都是常见字体.
  曲豆豆模板提供了这些字体的快捷输入法(自定义命令).
  部分自定义命令如下:
  
  \[
    \begin{tabular}{|c|c|c|c|}
    \hline
       字体
      &符号
      &通常命令
      &自定义命令
    \\ \hline
       空心体
      &$\bbR, \bbC, \bbA$
      & \verb"\mathbb{R}", \verb"\mathbb{C}" ,\verb"\mathbb{A}"
      & \verb"\bbR" ,\verb"\bbC" ,\verb"\bbA"
    \\ \hline
       哥特体
      & $\mfkp,\mfkm$
      & \verb"\mathfrak{p}",\verb"\mathfrak{m}"
      & \verb"\mfkp",\verb"\mfkm"
    \\ \hline
       手写体
      & $\mcalF, \mcalU$
      &  \verb"\mathcal{F}",\verb"\mathcal{U}"
      & \verb"\mcalF" , \verb"\mcalU"
    \\ \hline
       粗体
      & $\bfe, \bfv$
      & \verb"\boldsymbol{e}" , \verb"\boldsymbol{v}"
      & \verb"\bfe", \verb"\bfv"
    \\ \hline
    \end{tabular}
  \]
  
  模板文件\verb"QuDoudouSymbols.tex"
  当中记载了非常多的自定义命令,
  感兴趣者可自行查看,
  但是\textbf{不建议擅自增删更改}.
  如需新增自定义命令,请与曲豆豆沟通.
  
  \item 还支持其它常用的缩写,
  比如将\verb"\alpha"简写为\verb"\afa", 
  \verb"\gamma"简写为\verb"\gma"等等.
  特别推荐将\verb"\varphi"简写为\verb"\fai" , 
  将\verb"\varepsilon" 简写为\verb" \veps".
  亦详见模板文件\verb"QuDoudouSymbols.tex".
  
  \item 支持用\verb"\xymatrix"排版交换图表.
  具体用法自行百度.
  
  \item 支持用\verb"\begin{figure}[H]"
  来插入图片.
  \textbf{务必将图片放入\texttt{figures}文件夹中.}
\end{enumerate}

\section{交叉引用规范}

代数学千题解是一个多人协作的大项目,
打字员不可避免要交叉引用.
交叉引用需要注意规范,
以方便引用,被引用别人的工作.
假定读者熟悉\verb"\label",\verb"\ref"的用法.

\begin{enumerate}
  \item 请对每一个\verb"\section"
  手动添加标签\verb"\label{section-a-b}",
  例如$3.7$节与$15.2$节的标签分别为
  \verb"section-3-7"与\verb"section-15-2".
  
  \item 请为每一道习题添加标签
  \verb"\label{ex-a-b-c}",
  其中a,b,c分别为章号,节号,题号.
  例如12.1节的习题5的标签为 \verb"ex-12-1-5".
  
  \item 若大家严格依照以上两条,
  将极大方便大家引用章节,习题编号.
  只需要\verb"\ref{...}".
  
  \item 不支持对Dummit书中正文的引用.
  如遇到题目中引用此类,请大家设法绕开.
  
  \item 此外,代数学千题解还提供中英文术语对照索引.
  为完成这个索引,请大家在\textbf{内容概要}(以及部分习题)
  当中介绍新概念时(尤其涉及英文名词时),
  务必使用曲豆豆开发的含参命令\verb"\newdef{}{}".
  用此命令可自动生成一条中英对照条目.
  示例:
  \begin{center}
    \verb"若你不会用这个命令,则称曲豆豆是\newdef{猪}{pig}".
  \end{center}
  其编译效果如下:
  \begin{center}
    若你不会用这个命令,则称曲豆豆是\newdef{猪}{pig}.
  \end{center}
  更多示例见示例文本.
\end{enumerate}




\chapter{示例文本}

\section{Noether环与仿射代数集}
\label{section-15-1}

本章提到的环都默认是含幺交换环,且$1\neq 0$.

\subsection{内容提要}

在本节,记$\Bbbk$为给定的一个域.

\subsubsection{Noether环与$\Bbbk$-代数}

\begin{itemize}
  \item 我们在\ref{section-9-6}节已介绍Noether环的基本概念,
  以及关于Noether环的Hilbert基定理.
  Noether环有以下三个常见的等价定义:理想升链条件;
  任何理想都有限生成;
  任意一族理想都有(包含偏序)极大元.

  \item Noether环的商环、同态像仍为Noether环.
  此外,主理想整环一定是Noether环.

  \item 非Noether环的常见例子(供了解):
  无穷变元多项式环$\bbZ[x_1,x_2,x_3,...]$,
  闭区间$[0,1]$上的连续函数环$C[0,1]$等等...

  \item 如果存在非零环同态$\Bbbk\to R$,
  称环$R$为\textbf{$\Bbbk$-代数}.
  称$\Bbbk$-代数$R$是有限生成的,
  如果存在$R$的一组作为环的有限生成元.

  \item 对于$\Bbbk$-代数$R$,
  则$\Bbbk$自然视为$R$的一个子环.
  $R$也自然有$\Bbbk$-线性空间结构.
  $\Bbbk$上的多项式环$\Bbbk[x_1,x_2,...,x_n]$
  是有限生成$\Bbbk$-代数的典型例子.
  (但它作为$\Bbbk$-线性空间,不是有限维的.)
  事实上,任何有限生成$\Bbbk$-代数必同构于
  $\Bbbk[x_1,x_2,...,x_n]$的某个商环,
  其中$n$为某个(充分大的)正整数.
  有限生成$\Bbbk$-代数一定是Noether环.

  \item 设$R,S$为$\Bbbk$-代数,
  以及环同态$\fai:R\to S$.
  如果下图交换,则称$\fai:R\to S$为$\Bbbk$-代数同态:
  \[\xymatrix{
     R\ar[rr]^{\fai}
    &
    &S
  \\
    &\Bbbk
      \ar[ul]
      \ar[ur]
    &
  }\]
\end{itemize}

\subsubsection{仿射代数集}

给定域$\Bbbk$,记$\bbA^n:=\bbA_{\Bbbk}^n
:=\Bigset{(a_1,a_2,...,a_n)}{a_i\in\Bbbk}$,
称为$\Bbbk$上的\newdef{仿射$n$-空间}{affine $n$-space}.
称环$\Bbbk[\bbA^n]:=\Bbbk[x_1,x_2,...,x_n]$,
为$\bbA^n$的\newdef{坐标环}{coordinate ring}.

\begin{itemize}
  \item 想法来源于解多元高次代数方程组.
  注意到$\Bbbk[x_1,x_2,...,x_n]$当中的多项式,
  自然可以看作$\bbA^n$上的函数.
  对于$\Bbbk[x_1,x_2,...,x_n]$的非空子集$S$,
  则$S$自然视为一个方程组$f_i(x_1,x_2,...,x_n)=0,(\forall f_i\in S)$.
  这个方程组的解的全体(即 '解空间')自然看作$\bbA^n$的一个子集.
  我们记
  \[\mcalZ(S):=\Bigset{\bfa\in\bbA^n}{f(\bfa)=0,\,\forall f\in S}\]
  称为$S$的\newdef{割迹}{locus}(俗称$S$当中的多项式的'公共零点集').
  并且特别规定$\mcalZ(\vkong)=\bbA^n$.
  容易验证,对于任意$\vkong\neq S\subseteq\Bbbk[x_1,x_2,...,x_n]$,
  都有$\mcalZ(S)=\mcalZ(I)$,其中$I$是由$S$生成的理想.

  \item 对于$\bbA^n$的子集$V$,
  如果存在$\Bbbk[\bbA^n]$的理想$I$使得
  $V=\mcalZ(I)$,
  则称$V$为\newdef{仿射代数集}{affine algebraic set}.

  \item 设$f\in\Bbbk[x_1,x_2,...,x_n]$为非常值多项式,
  则$\mcalZ((f))\subseteq\bbA^n$为仿射代数集,
  这样的仿射代数集称为\newdef{超曲面}{hypersurface}.
  由Hilbert基定理易知任何仿射代数集都是有限多个超曲面之交.

  \item 对于仿射代数集$V\subseteq\bbA^n$,
  使得$\mcalZ(I)=V$的理想$I$显然不唯一:
  如果$\mcalZ(I)=V$,
  那么对任意$k\geq 1$都有$\mcalZ(I^k)=\mcalZ(I)=V$.
  于是我们企图去找满足此条件的"最大"的理想$I$,
  从而引入以下定义:
  对$\bbA^n$的任意子集$A$,定义
  \[
    \mcalI(A):=
    \Bigset{f\in\Bbbk[\bbA^n]}
           {f(\bfa)=0,\,\forall \bfa\in A}
  \]
  则$\mcalI(A)$是$\Bbbk[\bbA^n]$的理想.

  \item 给定$n\geq 0$,则有以下映射:
  \begin{eqnarray*}
    \mcalZ:
      \Big\{
        \text{$\Bbbk[\bbA^n]$的子集}
      \Big\}
    &\to&
    \Big\{
      \text{$\bbA^n$的仿射代数集}
    \Big\}
  \\
    \mcalI:
    \Big\{\text{$\bbA^n$的子集}\Big\}
    &\to&\Big\{\text{$\Bbbk[\bbA^n]$的理想}\Big\}
  \end{eqnarray*}

  \item 容易验证,对$\bbA^n$的任意子集$A$，
  以及$\Bbbk[\bbA^n]$的任意理想$I$,都成立
  \[\mcalZ(\mcalI(A))\supseteq A,\quad
    \mcalI(\mcalZ(I))\supseteq I\]
  一般来说上述包含关系未必取等.

  \item 但是,当$A$为仿射代数集时,
  当$I$形如$\mcalI(A)$时,可以证明必有
  \[\mcalZ(\mcalI(A))= A,\quad
    \mcalI(\mcalZ(I))= I,\]
  从而有如下一一对应:
  \[
    \xymatrix{
      \Big\{
        \bbA^n\text{的仿射代数集}
      \Big\}
      \ar@/^0.2pc/[r]^-\mcalI
    &
     \Bigset{\mcalI(A)}{A\subseteq\bbA^n}.
     \ar@/^0.2pc/[l]^-\mcalZ
    }
  \]
  不过要注意,并不是$\Bbbk[\bbA^n]$的所有理想都位于
  集合$\Bigset{\mcalI(A)}{A\subseteq\bbA^n}$当中,
  本节暂无对该集合的具体刻画.
\end{itemize}

\subsubsection{仿射代数集的态射及其坐标环同态}

本小节给定$m,n\geq 0$,
记$V\subseteq\bbA^m$,
$W\subseteq\bbA^n$为仿射代数集.

\begin{itemize}
\item 对于仿射代数集$V\subseteq\bbA^m$,
称$\Bbbk[V]:=\Bbbk[\bbA^m]/\mcalI(V)$
为$V$的\newdef{坐标环}{coordinate ring}.
$\Bbbk[V]$中的元素俗称'$V$上的多项式函数'.

\item 对于映射$\fai:V\to W$,
如果存在多项式$\fai_1,\fai_2,...,\fai_n\in\Bbbk[x_1,x_2,...,x_m]$,
使得对任意$\bfv\in V$都有$\fai(\bfv)=(\fai_1(\bfv),\fai_2(\bfv),...,\fai_n(\bfv))$,
则称$\fai:V\to W$为仿射代数集之间的\newdef{态射}{morphism},
也称为\newdef{代数映射}{algebriac map},
\newdef{正则映射}{regular map}.
对于代数映射$\fai:V\to W$,
如果存在代数映射$\psi:W\to V$使得$\fai\circ\psi=\id_W$
以及$\psi\circ\fai=\id_V$,
则称$\fai:V\to W$为仿射代数集$V,W$之间的一个
\newdef{同构}{isomorphism}.

\item 上述定义的"代数映射"无非是
"在局部坐标下能用多项式表示"的"多项式映射".
事实上代数映射还有以下内蕴的等价定义:
映射$f:V\to W$为代数映射,当且仅当对任意$f\in\Bbbk[W]$,
$f\circ\fai\in\Bbbk[V]$.

\item 仿射代数集之间的态射诱导其坐标环的$\Bbbk$-代数同态.
设$\fai:V\to W$为代数映射,
则定义$\Bbbk$-代数同态$\faitil:\Bbbk[W]\to\Bbbk[V]$如下:
对任意$f\in\Bbbk[W]$,将$f$视为$W$上的函数,
则$\faitil(f):=f\circ\fai\in\Bbbk[V]$.

\item 反之,对于任意$\Bbbk$-代数同态$\Phi:\Bbbk[W]\to\Bbbk[V]$,
必存在唯一的代数映射$\fai:V\to W$使得$\Phi=\faitil$.
具体地,记$\oly_1,\oly_2,...,\oly_n$
为$y_1,y_2,...,y_n\in\Bbbk[y_1,y_2,...,y_n]=\Bbbk[\bbA^n]$
在$\Bbbk[W]$当中的像,则对任意$\bfa\in V$,定义
\[\fai(\bfa):=(\Phi(\oly_1)(\bfa),\Phi(\oly_2)(\bfa),...,\Phi(\oly_n)(\bfa)),\]
则$\fai:V\to W$为代数映射,并且$\faitil=\Phi$.

\item 上述给出了一一对应:
\[
  \Big\{
    V\to W\,\text{代数映射}
  \Big\}
  \xle{1-1}
  \Big\{
    \Bbbk[W]\to\Bbbk[V],\text{$\Bbbk$-代数同态}
  \Big\}
\]
并且容易证明,设$\fai: V\to W$
以及$\psi: W\to U$为代数映射,则$\psi\circ\fai:W\to U$
也是代数映射,并且$\widetilde{\psi\circ\fai}=\faitil\circ\psitil:
\Bbbk[U]\to\Bbbk[V]$.
特别地,代数映射$\fai:V\to W$是仿射代数集之间的同构,
当且仅当$\faitil:\Bbbk[W]\to\Bbbk[V]$是$\Bbbk$-代数同构.
\end{itemize}

\subsubsection{$\Bbbk$-代数同态的具体计算}

设$I,J$分别为$\Bbbk[x_1,x_2,...,x_m]$,
$\Bbbk[y_1,y_2,...,y_n]$的理想,
\[
  \Phi:\Bbbk[y_1,y_2,...,y_n]/J
  \to
  \Bbbk[x_1,x_2,...,x_m]/I
\]
为任意给定的$\Bbbk$-代数同态.
我们可以利用Gr\"{o}bner基的理论(详见\ref{section-9-6}节)
来具体计算$\ker\Phi$,
并且判断$\Bbbk[x_1,x_2,...,x_m]$中的元素是否位于$\im\Phi$.
我们记
\[R:=\Bbbk[x_1,x_2,...,x_m; y_1,y_2,...,y_n],\]
并赋以$x_1>x_2>\cdots>x_m>y_1>y_2>\cdots>y_n$的字典序.
取定$\fai_1,\fai_2,...,\fai_n\in\Bbbk[x_1,x_2,...,x_m]$
使得$\fai_i$为$\Phi(\oly_i)$的一个代表元.
考虑$R$的理想
\[
  \mcalA:=(y_1-\fai_1,y_2-\fai_2,...,y_n-\fai_n)+I.
\]
记$G$为$R$的理想$\mcalA$的一组既约Gr\"{o}bner基.
\begin{itemize}
  \item 记号同上,则有$\ker\Phi=\mcalA\cap\Bbbk[y_1,y_2,...,y_n]$.
  则由消去定理(见\ref{section-9-6}节)可知$G\cap\Bbbk[y_1,y_2,...,y_n]$
  为$\ker\Phi\subseteq\Bbbk[y_1,y_2,...,y_n]$的一组Gr\"{o}bner基.
  \item 对于$f\in\Bbbk[x_1,x_2,...,x_m]$,
  则$\olf\in\im\Phi$当且仅当$f$做关于$G$的广义多项式除法所得的余式$h$
  位于$\Bbbk[y_1,y_2,...,y_m]$当中.
  并且此时成立$\Phi(\olh)=\olf$.
  \item 特别地,$\Phi$为满同态当且仅当对任意$1\leq j\leq n$,
  存在$f\in G$使得$\LT(f)=x_i$.
  (回顾:$\LT(f)$为多项式$f$的首项,详见\ref{section-9-6}节).

  \item 应用:设域扩张$\bbK/\bbF$为代数单扩张,
  $\bbK=\bbF(\afa)$,其中$\afa$为不可约多项式$p\in\bbF[x]$的一个根.
  则对于任意$f\in\bbF[x]$, $\beta:=f(\afa)\in\bbK$.
  我们可以通过以下方法来计算$\beta$在$\bbF$上的极小多项式:
  注意同构$\bbK\cong\bbF[x]/p(x)$,考虑$\bbF$-代数同态
  \begin{eqnarray*}
    \Phi:\bbF[x]&\to&\bbF[x]/p(x)\cong\bbK\\
              x&\mapsto&\overline{f(x)}
  \end{eqnarray*}
  则$\ker\Phi$的生成元即为$\beta$在$\bbF$上的极小多项式.
\end{itemize}

\subsection{习题与解答}

依然假定$\Bbbk$为给定的域,
所有提到的环都是含幺交换环并且$1\neq 0$.

\begin{enumerate}
  \begin{prob}\label{ex-15-1-1}
    证明Hilbert基定理的逆定理:
    若$R[x]$为Noether环,则环$R$也为Noether环.
  \end{prob}

  \begin{proof}
    若$R$不是Noethre环,
    取$R$的无限理想升链$I_1\subsetneqq I_2\subsetneqq I_3\subsetneqq\cdots$,
    则易知$(I_1x)\subsetneqq(I_2x)\subsetneqq(I_3x)\subsetneqq\cdots$
    为$R[x]$的理想升链,从而$R[x]$不是Noether环.
  \end{proof}

  \begin{prob}\label{ex-15-1-2}
    证明下述环不是Noether环,
    并且具体给出它们的一列无限理想升链:
    \begin{enumerate}
      \item 闭区间$[0,1]$上的实值连续函数环$C[0,1]$,
      \item 设$X$为无限集,考虑定义在$X$,
      取值于$\bbZ/2\bbZ$的函数构成的环$\Map(X,\bbZ/2\bbZ)$.
    \end{enumerate}
  \end{prob}

  \begin{proof}
    考虑对于$k=1,2,3,...$,
    $I_k:=\Bigset{f\in C[0,1]}{f(\frac{1}{j})=0,\,\forall j\geq k}$,
    则$I_k$为$C[0,1]$的理想,并且成立
    $I_1\subsetneqq I_2\subsetneqq I_3\subsetneqq \cdots$.
    再任取无限集$X$当中的一列互不相同的元素$x_1,x_2,...$,
    考虑$J_k:=\Bigset{g:X\to\bbZ/2\bbZ}{g(x_j)=0,\,\forall j\geq k}$,
    则$J_k$为环$\Map(X,\bbZ/2\bbZ)$的理想,并且
    $J_1\subsetneqq J_2\subsetneqq J_3\subsetneqq \cdots$.
  \end{proof}

  \begin{prob}\label{ex-15-1-3}
    证明:$\Bbbk$上的一元有理函数域$\Bbbk(x)$不是有限生成$\Bbbk$-代数.
    \hint{回顾$\Bbbk(x)$为多项式环$\Bbbk[x]$的分式环,
    域扩张$\Bbbk(x)/\Bbbk$是有限生成扩张.}
  \end{prob}

  \begin{proof}
    假设$\Bbbk(x)$是有限生成$\Bbbk$-代数,
    并且$\Bigset{\frac{P_i}{Q_i}}{1\leq i\leq m}$
    为其一组$\Bbbk$-代数生成元,
    其中$P_i,Q_i\in\Bbbk[x]$, $Q_i\neq 0$,
    并且$P_i$与$Q_i$互素.
    对每个$1\leq i\leq m$,
    设$Q_i=\prod_{j=1}^{n_i}p_{ij}^{d_{ij}}$
    为素因式分解,其中$p_{ij}$为$\Bbbk[x]$中的不可约多项式.
    由于$\Bigset{\frac{P_i}{Q_i}}{1\leq i\leq m}$
    为$\Bbbk(x)$的$\Bbbk$-代数生成元,
    从而对任意$\frac{P}{Q}\in\Bbbk[x]$,
    若$P,Q\in\Bbbk[x]$并且$P,Q$互素,
    则分母$Q$的任何素因子必形如$p_{ij}$,
    其中$1\leq i\leq m$, $1\leq j\leq n_i$.
    但另一方面,
    注意$\Bbbk[x]$中的不可约多项式有无限多个,
    任取$r\in\Bbbk[x]$为不可约多项式,
    且$r\not\in\Bigset{p_{ij}}{1\leq i\leq m,
    \, 1\leq j\leq n_i}$,
    则$\frac{1}{r}$无法由
    $\Bigset{\frac{P_i}{Q_i}}{1\leq i\leq m}$ $\Bbbk$-代数生成,
    故$\frac 1r\not\in\Bbbk(x)$,矛盾.
  \end{proof}

  \begin{prob}\label{ex-15-1-4}
    证明:若环$R$是Noether环,则形式幂级数环$R\fps{x}$
    (见\ref{section-7-2}节的习题\ref{ex-7-2-3})
    也是Noether环.
    \hint{模仿Hilbert基定理的证明.}
  \end{prob}

  \begin{proof}
    对于形式幂级数$0\neq f=\sum_{k=0}^{+\infty}a_kx^k\in R\fps{x}$,
    称$\deg f:=\max\Bigset{k}{a_0=a_1=\cdots=a_{k-1}=0,\, a_k\neq 0}$
    为$f$的阶数;
    若$\deg f=k$,
    则令$\LT(f)$,则称$a_k\in R$为$f$的首项系数.
    以下模仿Hilbert基定理的证明来证明本题:
    \begin{itemize}
      \item 设$I$为$R\fps{x}$的非零理想,
      只需证明$I$是$R\fps{x}$-有限生成的.考虑
      \[I_0:=\Bigset{a\in R}
      {\text{$a$为$I$当中某个形式幂级数的首项系数}}\cup\{0\},\]
      则易验证$I_0$为$R$的理想.
      由于$R$为Noether环,
      从而存在$a_1,a_2,...,a_m\in R$构成理想$I_0$的一组$R$-生成元.
      对每个$1\leq i\leq m$,
      取形式幂级数$f_i=x^{d_i}
      \sum_{k=0}^{+\infty}a_{ik}x^k\in R\fps{x}$,
      其中$a_{i0}=a_i\neq 0$,
      $\deg f_i=d_i\geq 0$.
      再记$d:=\max\Bigset{d_i}{1\leq i\leq m}$
      为$f_1,f_2,...,f_m$的最高阶数.

      \item 断言:$\{f_1,f_2,...,f_m\}$
      是理想$I\cap\Bigset{f\in I}{\deg f\geq d}$
      的一组$R\fps{x}$-生成元.
      这是因为,对任意$g\in I$并且$\deg g\geq d$,
      记$g_0:=g$.记$h_0:=\deg g_0$,
      $g_0$的首项系数为$b_0\in I_0\subseteq R$.
      从而存在$c_{0j}\in R,\,j=1,2,...,m$
      使得$b_0=\sum_{j=1}^{m}c_{0j}a_j$.

      考虑$g_1:=g_0-\sum_{j=1}^{m}c_{0j}x^{h_0-d_j}f_j$,
      则$g_1\in I$,并且$\deg g_1>\deg g_0=h_0\geq d$.
      记$h_1:=\deg g_1$,
      $g_1$的首项系数为$b_1\in I_0\subseteq R$.
      则存在$c_{1j}\in R,\,j=1,2,...,m$,
      使得$b_1=\sum_{j=1}^{m}c_{1j}a_j$.

      将上述操作反复进行下去,
      一般地,对于任意$n$,考虑
      \[g_{n}:=g_{n-1}-\sum_{j=1}^{m}c_{n-1,j}x^{h_{n-1}-d_j}f_j,\]
      则$g_n\in I$,并且$\deg g_n>\deg g_{n-1}>d$.
      记$h_n:=\deg g_n$,
      $g_n$的首项系数为$b_n\in I_0\subseteq R$.
      则存在$c_{nj}\in R,\, j=1,2,...,m$,
      使得$b_n=\sum_{j=1}^{n}c_{nj}a_j$.

      如此不断进行下去,会得到一列
      $\Bigset{c_{nj}}{1\leq j\leq m,\, n=0,1,2,...}$,
      以及一列$h_n,\,n\geq 0$.
      现在对任意$1\leq j\leq m$,令
      $\fai_j:=\sum_{n=0}^{+\infty}c_{nj}x^{h_n-d_j}\in R\fps{x}$.
      则显然有
      \[
        g=\sum_{j=1}^{m}\fai_jf_j
      \]
      而由$g$的任意性,我们证明了$I$中的任何阶数$\geq d$的形式幂级数
      都可被$\{f_1,f_2,...,f_m\}\,$ $R\fps{x}$-生成.

      \item 目前我们已保证$I$当中的阶数"较高"的形式幂级数能被
      $\{f_1,f_2,...,f_m\}$生成.
      我们只需再处理阶数"较低"的.
      对于每个$1\leq h<d$,考虑
      \[
        I_h:=\Bigset{a\in R}
        {\text{$a$为$I$中某个阶数为$h$的形式幂级数的首项系数}}
        \cup\{0\}
      \]
      则易验证$I_h$为$R$的理想.
      由于$R$为Noether环,
      从而存在$I_h$的一组有限生成元$\{a_{h1},a_{h2},...,a_{hm_h}\}\subseteq R$.
      对每个$a_{hj},\, 1\leq j\leq m_h$,
      取$f_{hj}\in R\fps{x}$,
      使得$\deg f_{hj}=h$,
      并且$f_{hj}$的首项为$a_{hj}$.

      \item 易证以下为$R\fps{x}$的理想$I$的一组有限生成元:
      \[
        \{f_1,f_2,...,f_m\}
      \cup
        \Bigset{f_{hj}}{1\leq h<d,\, 1\leq j\leq m_h}
      \]
      从而理想$I$为有限生成的,从而证毕.
      \remark{
       大致思路: 对任意$g\in I$,
       先用$\{f_{hj}\}$"杀掉"$g$的"低阶部分",
       再用$\{f_j\}$"杀掉"剩余的"高阶部分".
       细节留给读者.
      }
    \end{itemize}
  \end{proof}

  \begin{prob}\label{ex-15-1-5}
    \newdef{Fitting 引理}{Fitting's lemma}.
    设$M$为Noether $R$-模,$\fai:M\to M$为模自同态.
    证明:对于充分大的正整数$n$,
    有$\ker(\fai^n)\cap\im(\fai^n)=0$.
    并由此推出,如果$\fai$是满射,则$\fai$为$R$-模自同构.
    \hint{注意到$\ker\fai\subseteq\ker(\fai^2)
    \subseteq\ker(\fai^3)\subseteq\cdots$.}
  \end{prob}

  \begin{proof}
    注意到$M$的子模升链$\ker\fai\subseteq\ker(\fai^2)
    \subseteq\ker(\fai^3)\subseteq\cdots$.
    因为$M$为Noether $R$-模,
    故存在$N\geq 1$,
    使得$\ker(\fai^N)=\ker(\fai^{N+1})=\ker(\fai^{N+2})=\cdots$.
    则对任意$n\geq N$,
    以及$m\in\ker(\fai^n)$,
    则$\fai^n(m)=0$.
    如果还有$m\in\im(\fai^n)$,
    记$m=\fai^n(m')$,
    则$\fai^{2n}(m')=\fai^n(m)=0$,
    所以$m'\in\ker(\fai^{2n})=\ker(\fai^n)$,
    从而$m=\fai^n(m')=0$.
    这就证明了$\ker(\fai^n)\cap\in(\fai^n)=0$.

    如果$\fai:M\to M$为满同态,则对任意$n$都有$\im(\fai^n)=M$.
    从而Fitting引理迫使$\ker(\fai^n)=0$,其中$n$充分大.
    于是$\ker(\fai)\subseteq\ker(\fai^n)=0$,
    从而$\fai$也必为单同态,故$\fai$为$R$-模自同构.
  \end{proof}

  \begin{prob}\label{ex-15-1-6}
    设$0\to M'\to M\to M''\to 0$为$R$-模正合列.
    证明:$M$为Noether模当且仅当$M',M''$都为Noether模.
  \end{prob}

  \begin{proof}
    必要性见下题\remark{注意$M', M''$分别是$M$的子模,商模},
    只证充分性.
    假设$M',M''$都为Noether模,
    则对$M$的子模$N$,
    $(N+M')/M'$为$M''\cong M/M'$的子模,
    取其$R$-模生成元$\olf_1,\olf_2,...,\olf_r\in (N+M')/M'$,
    其中$f_j\in M$为$\olf_j$的一个代表元.
    再注意$N\cap M'$为$M$的子模,
    取其$R$-模生成元$g_1,g_2,...,g_s\in M$.
    则易知$\{f_1,f_2,...,f_r; g_1,g_2,...,g_s\}$
    为子模$N$的$R$-生成元.
    于是$M$的每个子模都是有限生成的,故$M$也为Noether模.
  \end{proof}

  \begin{prob}\label{ex-15-1-7}
    证明:Noether $R$-模的子模,商模,有限直和仍为Noether $R$-模.
  \end{prob}

  \begin{proof}
    设$M$为Noether $R$-模,
    $N$为$M$的子模.
    只需注意$N$的子模仍为$M$的子模,
    从而$N$的子模升链也是$M$的子模升链,
    从而必停止.
    再考虑$N$的商模$M/N$,
    只需注意一一对应$\{\text{$M/N$的子模}\}\xle{1-1}
    \{\text{$M$的包含$N$的子模}\}$.
    在此一一对应下,$M/N$的子模升链视为$M$的子模升链,
    从而必停止.故$M/N$也是Noether模.

    若$M',M''$都是Noether模,则注意到短正合列
    \[0\to M'\to M'\oplus M''\to M''\to 0 \]
    故由上题结论知$M'\oplus M''$也为Noether模.
  \end{proof}


  \begin{prob}\label{ex-15-1-8}
    设$R$为Noether环,则$R$-模$M$是Noether模
    当且仅当$M$是有限生成$R$-模.
    \remark{从而当$R$为Noether环时,
    有限生成$R$模的任何子模也是有限生成的.}
  \end{prob}

  \begin{proof}
    必要性显然,只证充分性.
    设$R$为Noether环,
    $M$为有限生成$R$-模,
    $\{m_1,m_2,...,m_s\}$为$M$的一组$R$-模生成元.
    注意$R$本身作为$R$-模也是Noether模,
    从而由上题可知$R^{\oplus s}$为Noether模.
    考虑$R$-模同态:
    \begin{eqnarray*}
      \fai:R^{\oplus s}&\surj& R\\
      e_i&\mapsto& m_i
    \end{eqnarray*}
  其中$e_1,e_2,...,e_s$为$R^{\oplus s}$的一组$R$-基.
  则$\fai$为满同态,并且有$R$-模同构$M\cong R^{\oplus s}/\ker\fai$.
  再利用上题,注意Noether模的商模仍为Noether模.
  \end{proof}

  \begin{prob}\label{ex-15-1-9}
    设$\Bbbk$为域,证明$\Bbbk[x]$的任何包含$\Bbbk$的子环都是Noether环.
    并且举例说明这样的环不一定是唯一分解整环.
    \hint{如果$\Bbbk\subseteq R\subseteq\Bbbk[x]$
    以及任意$y\in R\setminus\Bbbk$,
    证明$\Bbbk[x]$为有限生成$\Bbbk[y]$-模,
    再利用之前的习题.至于第二问,考虑$\Bbbk[x^2,x^3]$.}
  \end{prob}

  \begin{proof}
    $\Bbbk[x]$的包含$\Bbbk$的子环,
    无非就是$\Bbbk[x]$作为$\Bbbk$-代数的$\Bbbk$-子代数.
    设$R$为$\Bbbk[x]$的$\Bbbk$-子代数,
    事实上,断言$R$必为有限生成$\Bbbk$-代数,
    从而为Noether环.以下分若干步证明之.

    \begin{itemize}
      \item 对于任意$n\geq 0$,
      若$\Bigset{f\in R}{\deg f=n}$非空,
      则任取此集合中的一个首一多项式,记为$f_n$;
      若此集合为空,则记$f_n:=0$.
      则易知$\Bigset{f\in R}{\deg f\leq n}
      \subseteq\Bbbk[f_1,f_2,...,f_n],\,\forall n\geq 0$.

      \item 记$n_1:=\min\Bigset{n\geq 1}{f_n\neq 0}$,
      再令$d_1:=n_1=\deg f_{n_1}$.
      如果$\{f_{n_1}\}$是$R$的一族$\Bbbk$-代数生成元,
      则证明结束;
      否则$R\setminus\Bbbk[f_{n_1}]$非空,
      取首一多项式$f_{n_2}\in R\setminus\Bbbk[f_1,f_2,...,f_{n_1}]$,
      使得$n_2=\min\Bigset{\deg g}{f\in R\setminus\Bbbk[f_1,f_2,...f_{n_1}]}$.
      则必有$n_2$不被$d_1=n_1$整除,
      否则考虑多项式$g_2:=f_{n_2}-f_{n_1}^{\frac{n_2}{d_1}}$,
      则$\deg g_2<n_2$,故由$n_2$的定义可知$g_2\in\Bbbk[f_1,f_2,...,f_{n_1}]$,
      进而$f_{n_2}\in\Bbbk[f_{n_1}]$,矛盾.
      记$d_2$为$n_1,n_2$的最大公约数,
      则$d_1>d_2\geq 1$.

      \item 归纳地,
      假设我们已经构造出$n_1<n_2<\cdots<n_k$,
      使得$d_k$为$n_1,n_2,...,n_k$的最大公约数,
      使得$d_1>d_2>\cdots>d_k\geq 1$;
      并且对任意$2\leq l\leq k$,
      $f_{n_l}$是非空集合$R\setminus\Bbbk[f_{1},f_{2},...,f_{n_{l-1}}]$
      当中次数最低的首一多项式.
      则考虑$\{f_1,f_2,...,f_{n_k}\}$,
      注意$d_k$为$n_1,n_2,...,n_k$的最大公约数,
      则由初等数论可知存在充分大的正整数$M_k>n_k$,
      使得对任意$n\geq M_k$,
      都存在非负整数$r_{k1},r_{k2},...,r_{kk}$
      使得$nd_k=\sum_{j=1}^{k}r_{kj}d_j$.
      现在,假设$\{f_1,f_2,...,f_{d_kM_k}\}$
      依然不是$R$的$\Bbbk$-代数生成元,
      则取$f_{n_{k+1}}$为
      $R\setminus\Bbbk[f_1,f_2,...,f_{d_kM_k}]$
      当中的次数最低的首一多项式.
      断言$n_{k+1}$
      不被$d_k$整除,
      否则考虑多项式
      \[g_{k+1}:=f_{n_{k+1}}-\prod_{j=1}^{k}f_{n_j}^{r_{kj}}\in R\]
      则$\deg g_{k+1}<n_{k+1}$,
      从而由$n_{k+1}$的定义(次数最低性)可知
      $g_{k+1}\in\Bbbk[f_1,f_2,...,f_{d_kM_k}]$,
      进而$f_{n_{k+1}}\in\Bbbk[f_1,f_2,...,f_{d_kM_k}]$,
      矛盾.
      记$d_{k+1}$为$n_1,n_2,...,n_{n+1}$的最大公约数,
      则$d_{k}>d_{k+1}\geq 1$.
      归纳完毕.

    \item 上述归纳构造产生了一列$d_1>d_2>\cdots>d_k\geq 1$,
    故该归纳构造必终止于某个$k_0$,
    使得$d_{k_0}=1$,
    即$n_1,n_2,...,n_{k_0}$的最大公约数为$1$,
    故存在充分大的正整数$M\geq n_{k_0}$,
    使得对任意$n\geq M$,
    都存在非负整数$s_{1},s_2,...,s_{k_0}$
    使得$n=\sum_{j=1}^{k_0}s_jn_j$.
    取定此$M$,
    则显然$\{f_1,f_2,...,f_M\}$
    为$R$的一组$\Bbbk$-代数生成元,
    故$R$是有限生成$\Bbbk$-代数,
    从而为Noether环.
    \end{itemize}

    $\Bbbk[x]$的$\Bbbk$-子代数不一定是唯一分解整环,
    例如在$\Bbbk[x^2,x^3]$当中,
    \[x^6=x^2\cdot x^2\cdot x^2=x^3\cdot x^3\]
    而$x^2,x^3$都是$\Bbbk[x^2,x^3]$的不可约元.
  \end{proof}

  \begin{prob}\label{ex-15-1-10}
    证明:$\Bbbk[x,y]$的子环$R:=\Bbbk[x,x^2y,x^3y^2,...,x^iy^{i-1},...]$
    不是Noether环,从而不是有限生成$\Bbbk$-代数.
    \remark{这表明Noether环的子环未必是Noether环,
    有限生成$\Bbbk$-代数的子代数也未必是有限生成$\Bbbk$-代数.}
  \end{prob}

  \begin{proof}
    只需构造$R$的理想的一个无穷升链.
    对于任意$k\geq 1$,记$R$的理想
    $I_k:=(x,x^2y,...,x^ky^{k-1})$,
    则容易验证作为$\Bbbk$-线性空间,以下
    \[
    \Bigset{x^ry^s}{r\geq 2,\,0\leq s\leq r-2}
    \cup\{x,x^2y,...,x^ky^{k-1}\}\]
    为$I_k$的一组$\Bbbk$-基.
    故有理想升链$I_1\subsetneqq I_2\subsetneqq I_3\subsetneqq\cdots$,
    从而$R$不是Noether环.
  \end{proof}

  \begin{prob}\label{ex-15-1-11}
    设环$R$的每个素理想都是有限生成的.
    我们通过以下步骤来证明$R$是Noether环:
      \begin{enumerate}
        \item 如果$R$的非有限生成理想构成的集合非空,
        证明该集合当中必存在极大元,记为$I$,
        并且此时$R/I$是Noether环.
        \item 证明:存在$R$的有限生成理想$J_1,J_2$,
        使得它们包含$I$,并且满足$J_1J_2\subseteq I$,
        $J_1J_2$有限生成.
        \hint{注意$I$不是素理想.}
        \item 证明:$I/J_1J_2$是$J_1/J_1J_2$
        的有限生成$R/I$-子模.
        \hint{利用习题\ref{ex-15-1-8}.}
        \item 利用上一小问导出$I$是有限生成$R$-模,
        从而产生矛盾.从而证明$R$为Noether环.
      \end{enumerate}
  \end{prob}

  \begin{proof}
    易知环$R$的任何素理想都是有限生成的,
    我们要证明$R$的每一个理想都是有限生成的.
    \begin{enumerate}
      \item 记$\mcalS:=\Bigset{I\subseteq R}{\text{$I$为$R$的非有限生成理想}}$,
      则$\mcalS\neq\vkong$.
      对于$\mcalS$中的任何一个(包含偏序)链
      $\Bigset{I_t}{t\in\mcalT}$,
      考虑$I':=\bigcup_{t\in\mcalT}I_t$,
      则$I'$为$R$的理想.
      假设$I'$是有限生成的,
      $\{r_1,r_2,...,r_m\}$为$I'$的一组生成元.
      则存在$I_1,I_2,...,I_m\in\mcalS$
      使得$r_k\in I_k,\,1\leq k\leq n$.
      将$I_1,I_2,...,I_m$当中(包含偏序)最大者记为$I_0\in\mcalS$,
      则$\{r_1,r_2,...,r_m\}\subseteq I_0$,
      因为$I_0$非有限生成,故$(r_1,r_2,...,r_m)\subsetneqq I_0$,
      这就与$(r_1,r_2,...,r_m)=I'\supseteq I_0$矛盾.
      此矛盾表明$I$非有限生成,即$I'\in\mcalS$,
      进而$I$为链$\Bigset{I_t}{t\in\mcalT}$的一个上界.
      于是由Zorn引理,$\mcalS$之中必存在包含偏序极大元,记为$I$.

      对于$R$的理想$K\supseteq I$,
      $I$的极大性迫使$K$为有限生成理想,
      其生成元组在$R/I$中的像恰为$R/I$的理想$K/I$的一组生成元,
      从而$K/I$是$R/I$的有限生成理想.
      注意$R/I$的理想与$R$的包含$I$的理想之间的一一对应,
      从而$R/I$的任何理想都是有限生成的,故$R/I$为Noether环.

      \item 注意上述$I$非有限生成,从而由题设可知$I$不是素理想.
      从而存在$a_1,a_2\in R\setminus I$,
      使得$a_1a_2\in I$.
      令$J_1=I+(a_1)$, $J_2=I+(a_2)$为$R$的理想,
      则易知$J_1,J_2\supsetneqq I$,
      并且$J_1J_2\subseteq I$.
      而$I$在$\mcalS$中的极大性迫使$J_1,J_2$为有限生成理想.

      \item 注意$IJ_1\subseteq J_1J_2$,
      从而$R$的$R$-模结构自然诱导$J_1/J_1J_2$的$R/I$-模结构,
      并且$I/J_1J_2$是其$R/I$-子模.
      注意$J_1$为$R$的有限生成理想,
      容易验证其任意一组理想生成元在$J_1/J_1J_2$的像
      是$J_1/J_1J_2$的一组$R/I$-模生成元,
      故$J_1/J_1J_2$是有限生成$R/I$-模.
      注意$R/I$是Noether环,
      从而由习题\ref{ex-15-1-8}可知,
      其子模$I/J_1J_2$也是有限生成$R/I$-模.

      \item 设$\{\olb_1,\olb_2,...,\olb_r\}$
      为$I/J_1J_2$的一组$R/I$-模生成元,
      其中$b_1,b_2,...,b_r\in I$使得$b_i$在$I/J_1J_2$的像为$\olb_i$.
      注意$J_1,J_2$都是$R$的有限生成理想,
      从而乘积理想$J_1J_2$也有限生成,
      记$\{c_1,c_2,...,c_s\}\subseteq J_1J_2\subseteq I$
      为理想$J_1J_2$的一组生成元.
      则对任意$f\in I$,
      存在$\lmd_1,\lmd_2,...,\lmd_r\in R$,
      使得在$R/I$-模$I/J_1J_2$当中成立
      \[\olf=\sum_{i=1}^{r}\overline{\lmd}_i\olb_i\]
      从而在$R$当中成立$f-\sum_{i=1}^{r}\lmd_ib_i\in J_1J_2$.
      因此理想$I$可由$\{b_1,b_2,...,b_r\}\cup\{c_1,c_2,...,c_s\}$有限生成.
      这与$I$非有限生成矛盾.
      导致此矛盾的原因是最初假定$\mcalS\neq\vkong$.
      因此必有$\mcalS=\vkong$,
      即$R$的任何理想都是有限生成的,得证.

    \end{enumerate}
  \end{proof}

  \begin{prob}\label{ex-15-1-12}
    设$R$为Noether环,$S$为有限生成$R$-代数.
    设$T\subseteq S$为$R$-代数,
    并且$S$为有限生成$T$-模,
    证明: $T$为有限生成$R$-代数.
    \hint{
      若$\{s_1,s_2,...,s_n\}$为$S$的一组$R$-代数生成元,
      $\{s_1',s_2',...,s_m'\}$为$S$的一组$T$-模生成元.
      证明$s_i,s_j's_k'$可被表示为$\{s_i'\}$的(有限)$T$-线性组合.
      若$T_0$为由这些$T$-线性组合系数生成的$R$-子代数,
      证明$S$作为$T_0$-模可被$\{s_i'\}$有限生成,
      由此推出$T$为有限生成$R$-代数.
    }
  \end{prob}

  \begin{proof}\,

  \begin{itemize}
    \item 由题设,取$\{s_1,s_2,...,s_n\}$
    为$S$的一组$R$-代数生成元,
    $\{s_1',s_2',...,s_m'\}$
    为$S$的一组$T$-模生成元.
    则对于任意$1\leq i\leq n$,
    存在$t_{ij}\in T,\, 1\leq j\leq m$,使得
    \[s_i=\sum_{j=1}^{m}t_{ij}s_j'\]
    对任意$1\leq\afa,\beta\leq m$,
    存在$u_{\afa\beta k}\in T,\,1\leq k\leq m$,使得
    \[s_{\afa}'s_{\beta}'=\sum_{k=1}^{m}u_{\afa\beta k}s_k'.\]

    \item 考虑由$\Bigset{t_{ij}}{1\leq i\leq n,\,1\leq j\leq m}\cup
    \Bigset{u_{\afa\beta k}}{1\leq\afa,\beta,k\leq m}$
    生成的$R$-代数,记为$T_0$.
    则$R\subseteq T_0\subseteq T\subseteq S$.
    注意$R$为Noether环,并且$T_0$为有限生成$R$-代数,
    从而由Hilbert基定理可知$T_0$为Noether环.

    \item 断言: $\{s_1',s_2',...,s_m'\}$
    是$S$的一组$T_0$-模生成元.
    这是因为,对任意$f\in S$,
    不妨$f=s_{i_1}s_{i_2}\cdots s_{i_k}$,
    则有(按Einstein求和约定对相同指标求和)
    \begin{eqnarray*}
       f
    &=&
      s_1',s_2',...,s_m'
    =
      (t_{i_1j_1}s_{j_1}')
      (t_{i_2j_2}s_{j_2}')
      \cdots
      (t_{i_kj_k}s_{j_k}')\\
    &=&
      (t_{i_1j_1}t_{i_2j_2}\cdots t_{i_kj_k})
      s_{j_1}'s_{j_2}'\cdots s_{j_k}'\\
    &=&
      (t_{i_1j_1}t_{i_2j_2}\cdots t_{i_kj_k})
      (u_{j_1j_2l_2}s_{l_2}')s_{j_2}'\cdots s_{j_k}'\\
    &=&
      (t_{i_1j_1}t_{i_2j_2}\cdots t_{i_kj_k})
      (u_{j_1j_2l_2}u_{l_2j_2l_3})s_{l_3}'s_{j_3}'\cdots s_{j_k}'\\
    &=&\cdots\\
    &=&
      \underbrace{(t_{i_1j_1}t_{i_2j_2}\cdots t_{i_kj_k})
      (u_{j_1j_2l_2}u_{l_2j_2l_3}\cdots u_{l_{k-1}j_{k-1}l_k})}_{\in T_0}
      s_{l_k}'
    \end{eqnarray*}
    这就证明了$S$为有限生成$T_0$-模,
    且$\{s_1',s_2',...,s_m'\}$为其一组$T_0$-模生成元.

    \item 注意$T_0$为Noether环,
    从而有限生成$T_0$-模$S$子模$T$为有限生成$T_0$-模.
    (见习题\ref{ex-15-1-8}).
    设$\{x_1,x_2,...,x_{m'}\}$为$T$的一组$T_0$-模生成元,
    再注意$\{t_{ij},u_{\afa\beta k}\}$为$T_0$的$R$-代数生成元,
    从而容易验证$\{t_{ij},u_{\afa\beta k},x_l\}$
    为$T$的$R$-代数生成元,其中$1\leq i\leq n,\,1\leq j,\afa,\beta,k\leq m$,
    以及$1\leq l\leq m'$.
    因此$T$为有限生成$R$-代数.
  \end{itemize}
  \end{proof}

  \begin{prob}\label{ex-15-1-13}
   设$S,T$为$\Bbbk[\bbA^n]$的子集,
   $A,B$为$\bbA^n$的子集,
   验证映射$\mcalZ$与$\mcalI$的如下基本性质:
   \begin{enumerate}
     \item 若$S\subseteq T$,则$\mcalZ(T)\subseteq\mcalZ(S)$.
     \item $\mcalZ(S)=\mcalZ(I)$,其中$I$为$S$的生成的$\Bbbk[\bbA^n]$的理想.
     \item 设$\Bigset{S_i}{i\in\mcalA}$为$\Bbbk[\bbA^n]$的任意一族子集,
     则成立$\bigcap_{i\in\mcalA}\mcalZ(S_i)=\mcalZ(\bigcup_{i\in\mcalA}S_i)$.
     \item $\mcalZ(I)\cup\mcalZ(J)=\mcalZ(IJ)$.
     \item $\mcalZ(0)=\bbA^n$, $\mcalZ(1)=\vkong$.
     \item 若$A\subseteq B$,则$\mcalI(B)\subseteq\mcalI(A)$.
     \item $\mcalI(A\cup B)=\mcalI(A)\cap\mcalI(B)$.
     \item $\mcalI(\vkong)=\Bbbk[\bbA^n]$.
     如果$\Bbbk$不是有限域,则$\mcalI(\bbA^n)=0$.
     \item $A\subseteq\mcalZ(\mcalI(A))$,
           $I\subseteq\mcalI(\mcalZ(I))$.
     \item $\mcalZ(I)=\mcalZ(\mcalI(\mcalZ(I)))$,
           $\mcalI(A)=\mcalI(\mcalZ(\mcalI(A)))$.
   \end{enumerate}
  \end{prob}

  \begin{proof}
    \begin{itemize}
      (a)-(g)是显然的,直接定义验证即可.
      至于(h), $\mcalI(\vkong)=\Bbbk[\bbA^n]$显然,
      而当$\Bbbk$不是有限域时,
      先证明$\mcalI(\bbA^1)=0$.
      这是因为对任意$f\in\mcalI(\bbA^1)$,
      则任意$x\in\bbA^1=\Bbbk$都是多项式$f\in\Bbbk[x]$的根.
      而$\Bbbk$不是有限域,从而$f$有无穷多个根,
      这迫使$f=0$.从而$\mcalI(\bbA^1)=0$.
      现在假设$\mcalI(\bbA^n)=0$,
      则对任意$f\in\mcalI(\bbA^{n+1})\subseteq\Bbbk[x_1,x_2,...,x_n; y]$,
      记$f(\bfx,y)=\sum_{k=1}^{m}b_k(\bfx)y^k$,
      其中$b_k\in\Bbbk[\bfx]=\Bbbk[x_1,x_2,...,x_n]$.
      对任意的$\bfx\in\bbA^n$,
      视$f$为关于$y$的一元多项式,
      由于对任意$y\in\bbA^1=\Bbbk$,$f(\bfx,y)=0$,
      从而迫使各项系数$b_k(\bfx)$均为$0$.
      再由$\bfx\in\bbA^n$的任意性,
      可知$n$元多项式$b_1,b_2,...,b_m\in\mcalI(\bbA^n)$,
      故由归纳假设知它们都为$0$,
      因此$f=0$.
      这就归纳证明了当$\Bbbk$不是有限域的时候$\mcalI(\bbA^n)=0$.

      (i)是显然的,容易验证.
      至于(j),利用(i)以及(a),(f)也容易证明:
      注意$I\subseteq \mcalI(\mcalZ(I))$,
      两边取$\mcalZ$得到$\mcalZ(I)\supseteq\mcalZ(\mcalI(\mcalZ(I)))$;
      而另一方面,将$\mcalZ(\mcalI(A))\supseteq A$
      当中的$A$换成$\mcalZ(I)$,
      就有$\mcalZ(I)\subseteq\mcalZ(\mcalI(\mcalZ(I)))$,
      于是综上所述就有$\mcalZ(I)=\mcalZ(\mcalI(\mcalZ(I)))$.
      而另一式也同理.
    \end{itemize}
  \end{proof}

  \begin{prob}\label{ex-15-1-14}
    对任何域$\Bbbk$,
    证明: $\bbA^1$的仿射代数集必形如$\vkong,\Bbbk$或者$\Bbbk$的有限子集.
  \end{prob}

  \begin{proof}
    设$I$为$\Bbbk[x]$的理想.
    若$I=0$,则$\mcalZ(I)=\bbA^1$;
    若$I=\Bbbk[x]$,则$\mcalZ(I)=\vkong$.
    而当$I$时$\Bbbk[x]$的非零真理想时,
    注意$\Bbbk[x]$为主理想整环,
    从而存在非零多项式$f\in\Bbbk[x]$使得$I=(f)$.
    此时$\mcalZ(I)$为多项式$f$的零点集,
    从而为有限集.
  \end{proof}

  \begin{prob}\label{ex-15-1-15}
    设$\Bbbk=\bbF_2$,
    $V=\{(0,0),(1,1)\}\subseteq\bbA^2$.
    证明: $\mcalI(V)=\mfkm_1\mfkm_2$,其中
    $\mfkm_1=(x,y)$, $\mfkm_2=(x-1,y-1)$.
  \end{prob}

  \begin{proof}
    显然$\mfkm_1\mfkm_2\subseteq\mcalI(V)$.
    另一方面,
    对任意$f(x,y)\in\mcalI(V)$,
    注意$x(x-1)=x^2-x\in\mfkm_1\mfkm_2$,
    以及$y(y-1)=y^2-y\in\mfkm_1\mfkm_2$.
    对$f(x,y)$作关于多项式$\{x^2-x,y^2-y\}$
    的关于字典序$x>y$的广义多项式除法(见\ref{section-9-6}节),
    记其余式为$g(x,y)$.
    则$f(x,y)-g(x,y)\in\mfkm_1\mfkm_2\subseteq\mcalI(V)$,并且
    $g(x,y)$必形如$axy+bx+cy+d$,其中$a,b,c,d\in\bbF_2$.
    注意$g(0,0)=g(1,1)=0$,
    从而$d=0$, $a+b+c=0$.
    因此$g(x)$无非以下四种情况,分类讨论之:
    \begin{eqnarray*}
      0    &=& 0\in\mfkm_1\mfkm_2\\
      xy+x &=& x(y-1)\in\mfkm_1\mfkm_2\\
      xy+y &=& y(x-1)\in\mfkm_1\mfkm_2\\
      x+y  &=& x(y-1)+y(x-1)\in\mfkm_1\mfkm_2
    \end{eqnarray*}
    总之必有$g(x,y)\in\mfkm_1\mfkm_2$.
    因此$f(x,y)\in\mfkm_1\mfkm_2$.
    这就证明了$\mcalI(V)=\mfkm_1\mfkm_2$.
  \end{proof}

  \begin{prob}\label{ex-15-1-16}
    设$V$为$\bbA^n$的仿射代数集,且为有限集.
    若$V$当中有$m$个元素,证明:
    $\Bbbk[V]$作为$\Bbbk$-代数,同构于$\Bbbk^m$.
    \hint{利用中国剩余定理.}
  \end{prob}

  \begin{proof}
    设$V=\{\bfp_1,\bfp_2,...,\bfp_m\}\subseteq\bbA^n$.
    考虑$\Bbbk$-代数同态
    \begin{eqnarray*}
      \fai:\Bbbk[\bfx]&\to&\Bbbk^m\\
      f&\mapsto&(f(\bfp_1),f(\bfp_2),...,f(\bfp_n))
    \end{eqnarray*}
    其中$\bfx=(x_1,x_2,...,x_n)$.
    则显然$\ker\fai=\mcalI(V)$.
    再断言$\fai$为满同态.
    对每个$\bfp_k=(p_{k1},p_{k2},...,p_{kn})\in V$,
    考虑$\Bbbk[\bfx]$的理想
    $\mfkm_k:=(x_1-p_{k1}, x_2-p_{k2},...,x_n-p_{kn})$.
    则$\Bigset{\mfkm_k}{1\leq k\leq m}$
    为$\Bbbk[\bfx]$的两两互异的极大理想,
    从而两两互素.
    故由中国剩余定理,商映射$\Bbbk[x]\surj\Bbbk[x]/\mfkm_k,\,1\leq k\leq m$
    诱导的环同态
    \begin{eqnarray*}
      \fai':\Bbbk[x]&\to&\prod_{k=1}^{m}\Bbbk[x]/\mfkm_k
    \end{eqnarray*}
    为满同态.
    再注意到对每个$1\leq k\leq m$,
    环同构$\Bbbk[\bfx]/\mfkm_k\cong\Bbbk$由
    $\olf\mapsto f(\bfp_k)$所诱导.
    从而可知$\fai$为满同态.
    因此$\Bbbk^n\cong\Bbbk[x]/\ker\fai=\Bbbk[x]/\mcalI(V)=\Bbbk[V]$.

  \end{proof}

  \begin{prob}\label{ex-15-1-17}
    设$\Bbbk$为有限域.
    证明: $\bbA^n$的任何子集都是仿射代数集.
  \end{prob}

  \begin{proof}
    对于一般的域$\Bbbk$,
    众所周知$\bbA^n$的任何有限子集都是仿射代数集.
    而当$\Bbbk$为有限域时,
    $\bbA^n$是有限集,
    从而$\bbA^n$的任何子集也是有限集,
    故为仿射代数集.
  \end{proof}

  \begin{prob}\label{ex-15-1-18}
    设$\Bbbk=\bbF_q$为$q$元有限域.
    证明: $\mcalI(\bbA^1)=(x^q-x)\subseteq\Bbbk[x]$.
  \end{prob}

  \begin{proof}
    由有限域的相关知识,
    对任意$a\in\bbF_q$都有$a^q=a$,
    从而$(x^q-x)\subseteq\mcalI(\bbA^1)$.
    另一方面,设$f\in\mcalI(\bbA^1)$,
    则任意$a\in\bbF_q$都有$f(a)=0$,
    从而在$\bbF_q(x)$当中成立$(x-a)|f(x)$.
    由$a$的任意性,以及$\Bigset{x-a}{a\in\bbF_q[x]}$
    为两两不同的不可约多项式,从而必有
    $\prod_{a\in\bbF_q}(x-a)$整除$f$.
    再注意$\bbF_q[x]$上的恒等式
    $\prod_{a\in\bbF_q}(x-a)=x^q-x$,
    立刻得到$f\in(x^q-x)$,
    故$\mcalI(\bbA^1)\subseteq(x^q-x)$.
    综上, $\mcalI(\bbA^1)=(x^q-x)$.
  \end{proof}

  \begin{prob}\label{ex-15-1-19}
    对任意非常值多项式$f\in\Bbbk[x]$,
    试利用$\Bbbk[x]$的唯一因式分解定理来描述$\mcalZ(f)$,
    进而描述$\mcalI(\mcalZ(f))$.
    证明: $\mcalI(\mcalZ(f))=(f)$
    当且仅当$f$可被分解为互异的一次因式的乘积.
  \end{prob}

  \begin{proof}
    设$f\in\Bbbk[x]$,
    考虑素因子分解$f(x)=\prod_{k=1}^{m}p_k(x)^{d_k}$,
    其中$\{p_k(x)\}$为$\Bbbk[x]$当中两两互异的不可约多项式,
    $d_1,d_2,...,d_m$为正整数.
    记$\mcalZ(f)=\{a_1,a_2,...,a_n\}$.
    则注意到
    \[g(x):=\prod_{1\leq k\leq m\atop \deg p_k=1}p_k(x)
    =\prod_{j=1}^{n}(x-a_j),\]
    易知$\mcalZ(f)=\mcalZ(g)$,
    从而$(g)\subseteq\mcalI(\mcalZ(f))$.
    另一方面,任意$h\in\mcalI(\mcalZ(f))
    =\mcalI(\{a_1,a_2,...,a_n\})$,
    则$a_1,a_2,...,a_n$都为$h$的根,
    从而$h$被$(x-a_1)(x-a_2)\cdots(x-a_n)=g(x)$整除,
    因此$h\in(g)$,从而$\mcalI(\mcalZ(f))\subseteq(g)$.
    综上, $\mcalI(\mcalZ(f))=(g)$.
    也就是说, $\mcalI(\mcalZ(f))$是由$f$的所有不同一次因式的乘积生成的理想.
    特别地, $\mcalI(\mcalZ(f))=(f)$当且仅当$f$为若干互异的一次因子的乘积.
  \end{proof}

  \begin{prob}\label{ex-15-1-20}
    设$f,g\in\Bbbk[x,y]$为不可约多项式,
    并且它们不相伴(即$f$不被$g$整除, $g$也不被$f$整除).
    证明: $\mcalZ(f,g)$为空集或者$\bbA^2$的有限子集.
    \hint{若$(f,g)\neq (1)$,
    试证明$(f,g)$中存在$\Bbbk[x]$中的非零多项式,
    同理也存在$\Bbbk[y]$中的非零多项式.
    记$R:=\Bbbk[x]$,$F=\Bbbk(x)$,
    利用Gauss引理证明$f,g$在$F[y]$当中互素.}
  \end{prob}

  \begin{proof}
    记$R=\Bbbk[x]$,$F=\Bbbk(x)$为$R$的分式域.
    将$f,g$视为$R=\Bbbk[x]$上的关于变元$y$的多项式.
    $f,g$在$\Bbbk[x,y]$不可约,
    从而在$R[y]$上也不可约.
    注意$R$是唯一分解整环,
    从而由Gauss引理知它们在$F[y]$上也都不可约.
    又由于$f,g$不相伴,故$f,g$在$F[y]$互素.
    从而存在$r,s\in F[y]$使得$rf+sg=1$.
    取多项式$M(x)\in\Bbbk[x]$
    使得$M(x)r(x,y)$与$M(x)s(x,y)$都属于$\Bbbk[x,y]$.则
    \[M(x)=M(x)\cdot 1=M(x)r(x,y)f(x,y)+M(x)s(x,y)g(x,y)\in(f,g)\]
    这就证明了理想$(f,g)\subseteq\Bbbk[x,y]$
    当中存在关于$x$的一元多项式$M(x)$.
    同理,$(f,g)$当中也存在关于$y$的一元多项式$N(y)$.从而
    \[\mcalZ(f,g)\subseteq
    \Bigset{(u,v)\in\bbA^2}{M(u)=N(v)=0}\]
    而上式右边为有限集,
    因此$\mcalZ(f,g)$为有限集(或空集).
  \end{proof}

  \begin{prob}\label{ex-15-1-21}
    将$2\times 2$的$\Bbbk$-矩阵
    $\begin{pmatrix}
       a & b \\
       c & d
     \end{pmatrix}$
    视为$\bbA^4$当中的点$(a,b,c,d)$.
    证明: $\SL_2(\Bbbk)$为$\bbA^4$的仿射代数集.
  \end{prob}

  \begin{proof}
    $\SL_2(\Bbbk)=\Bigset{A\in\Bbbk^{2\times 2}}{\det A=1}$,
    视为$\bbA^4$的子集,无非是
    $\Bigset{(a,b,c,d)\in\bbA^4}{ad-bc-1=0}$,
    这里$ad-bc-1$视为$\Bbbk[a,b,c,d]$当中的多项式.
    从而$\SL_2(\Bbbk)$为$\bbA^4$的仿射代数集.
  \end{proof}

  \begin{prob}\label{ex-15-1-22}
    证明: $\SL_n(\Bbbk)$为$\bbA^{n^2}$的仿射代数集.
    \hint{这是上题结论的推广.}
  \end{prob}

  \begin{proof}
    记$\{x_{ij}\},\,1\leq i,j\leq n$为$n^2$的个自由变元.
    则$\det:=\det\begin{pmatrix}
                  x_{11} & x_{12} & \cdots & x_{1n} \\
                  x_{21} & x_{22} & \cdots & x_{2n} \\
                  \vdots & \vdots & \ddots & \vdots \\
                  x_{n1} & x_{n2} & \cdots & x_{nn}
                \end{pmatrix}$
  为关于这$n^2$个自由变元的$n$次多项式.
  将$\SL_n(\Bbbk)$视为$\bbA^{n^2}$的子集,
  有$\SL_n(\Bbbk)=\mcalZ(\det-1)$为仿射代数集.
  \end{proof}

  \begin{prob}\label{ex-15-1-23}
    设$V$为$\bbR^2$当中的直线
    (即某个非零的一次多项式$ax+by+c$的零点集),
    证明: 坐标环$\bbR[V]$作为$\bbR$-代数,
    同构于$\bbR[x]$,
    并给出$V$与$R$的仿射代数集的同构.
  \end{prob}

  \begin{solution}
    不妨$a\neq 0$ ($b\neq 0$的情形类似).
    记$t$为$\bbA^1$的坐标,$(x,y)$为$\bbA^2$的坐标.
    考虑代数映射
    \begin{eqnarray*}
      \fai:\bbA^1&\to& V\\
       t&\mapsto&(-\frac bat-\frac ca, t)\\
      \psi:V&\to& \bbA^1\\
       (x,y)&\mapsto& y
    \end{eqnarray*}
    容易验证$\fai$与$\psi$互为逆映射,
    故有仿射代数集的同构$V\cong\bbA^1$,
    从而有$\bbR$-代数同构$\bbR[V]\cong\bbR[x]$.
  \end{solution}

  \begin{prob}\label{ex-15-1-24}
    记$V=\mcalZ(xy-z)\subseteq\bbA^3$.
    证明: 仿射代数集$V$同构于$\bbA^2$.
    并具体给出同构映射$\fai:V\to\bbA^2$,
    $\faitil:\Bbbk[\bbA^2]\to\Bbbk[V]$,
    以及它们的逆映射.
    仿射代数集$\mcalZ(xy-z^2)$与$\bbA^2$同构吗?
  \end{prob}

  \begin{proof}\,
    \begin{itemize}
    \item
    记$\bbA^3$的坐标为$(x,y,z)$,
    $\bbA^2$的坐标为$(u,v)$.
    注意$V=\Bigset{(x,y,xy)\in\bbA^3}{x,y\in\Bbbk}$.
    考虑代数映射
    \begin{eqnarray*}
      \psi: \bbA^2&\to& V\\
      (u,v)&\mapsto&(u,v,uv)\\
      \fai:V&\to&\bbA^2\\
      (x,y,xy)&\mapsto& (x,y)
    \end{eqnarray*}
  则易知$\fai,\psi$互为逆映射,
  从而它们诱导了仿射代数集之间的同构$V\cong\bbA^2$.
  于是,$\faitil:\Bbbk[\bbA^2]\to\Bbbk[V]$
  为$\Bbbk$代数同构,
  且满足$\faitil(u)=\olx,\,\faitil(v)=\oly$.
  $\faitil^{-1}=\psitil:\Bbbk[V]\to\Bbbk[\bbA^2]$满足
  $\begin{cases}
    \psitil(\olx)=u\\
    \psitil(\oly)=v\\
    \psitil(\olz)=uv
  \end{cases}.$

  \item 至于$\mcalZ(xy-z^2)$与$\bbA^2$是否同构,
  这与域$\Bbbk$的选取有关.
  例如当$\Bbbk=\bbF_2$时,
  $\mcalZ(xy-z^2)=\{(0,0,0),(0,1,0),(1,0,0),(1,1,1)\}$
  共有四个点,
  从而由习题$\ref{ex-15-1-16}$可知
  $\Bbbk[\mcalZ(xy-z^2)]\cong\Bbbk^4$.
  同理$\Bbbk[\bbA^2]\cong\Bbbk^4$,
  故$\Bbbk[\bbA^2]\cong\Bbbk[\mcalZ(xy-z^2)]$,
  所以有仿射代数集的同构$\mcalZ(xy-z^2)\cong\bbA^2$.
  然而当$\Bbbk=\bbR$时,
  断言$\mcalZ(xy-z^2)$与$\bbA^2$不同构.
  若$\fai:\bbA^2\to\mcalZ(xy-z^2)$为同构映射,
  取$\bbA^3,\bbA^2$通常的欧氏拓扑,
  则$\fai$作为多项式映射必然连续,$\fai^{-1}$也连续,
  从而$\bbA^2$与$\mcalZ(xy-z^2)$(在欧氏拓扑下)同胚.
  但是,$\mcalZ(xy-z^2)$去掉点$(0,0,0)$之后有两个连通分支,
  而$\bbA^2$去掉任何一点后都任然连通,
  这表明$\mcalZ(xy-z^2)$与$\bbA^2$不同胚,矛盾.
  \end{itemize}
  \end{proof}

  \begin{prob}\label{ex-15-1-25}
    设$V\subseteq\bbA^n$为仿射代数集,
    $f\in\Bbbk[V]$.
    $f$的图像$\Gma(f):=
    \Bigset{(\bfa,b)\in\bbA^{n+1}}{\bfa\in V,\,b=f(\bfa)}$.
    证明: $\Gma(f)$为$\bbA^{n+1}$的仿射代数集,
    并且同构于$V$.
    \hint{考虑$\fai:V\to\Gma(f)$,使得$\fai(\bfa)=(\bfa,f(\bfa))$.}
  \end{prob}

  \begin{proof}
    设$V=\mcalZ(I)$,
    其中$I$为$\Bbbk[\bbA^n]$的理想.
    记$\bbA^{n+1}$的坐标为$(\bfx,y)$,
    其中$\bfx=(x_1,x_2,...,x_n)$为$\bbA^n$的坐标.
    任取$f\in\Bbbk[V]$在$\Bbbk[\bbA^n]$中的一个代表元,
    仍记为$f$.
    则易知$\Gma(f)=\mcalZ(I+(y-f))\subseteq\bbA^{n+1}$.
    故$\Gma(f)$为仿射代数集.
    考虑代数映射
    \begin{eqnarray*}
      \fai: V&\to&\Gma(f)\\
      \bfx&\mapsto&(\bfx,f(\bfx))\\
      \psi:\Gma(f)&\to& V\\
      (\bfx,y)&\mapsto&\bfx
    \end{eqnarray*}
  则易知$\fai,\psi$互为逆映射,
  因此有仿射代数集的同构$V\cong\Gma(f)$.
  \end{proof}

  \begin{prob}\label{ex-15-1-26}
    设$V=\mcalZ(xz-y^2,yz-x^3,z^2-x^2y)\subseteq\bbA^3$.
    \begin{enumerate}
      \item 证明: $\fai:\bbA^1\to V$, $t\mapsto(t^3,t^4,t^5)$为满态射.
      \hint{对于$(x,y,z)\neq(0,0,0)$,考虑$t=\frac yx$.}
      \item 试具体描述$\faitil:\Bbbk[V]\to\Bbbk[\bbA^1]$.
      \item 证明: $\fai$不是仿射代数集之间的同构.
    \end{enumerate}
  \end{prob}

  \begin{proof}
    容易验证对任意$t\in\bbA^1$,
    $\fai(t)=(t^3,t^4,t^5)\in\mcalZ(xz-y^2,yz-x^3,z^2-x^2y)$,
    从而$\fai$的确是从$\bbA^1$到$V$的态射.

    \begin{enumerate}
      \item 对任意$(x,y,z)\in V$,
      则有$
      \begin{cases}
        xz=y^2\\
        yz=x^3\\
        x^2y=z^2
      \end{cases}
      $,
      不妨$z\neq 0$,则$x,y$也非零.从而易知$x=(\frac yx)^3$
      , $y=(\frac yx)^4$以及$z=(\frac yx)^5$,
      因此考虑$\frac yx\in\bbA^1$,
      则有$\fai(\frac yx)=(x,y,z)$.
      因此$\fai$为满态射.
      \item $\faitil$将$\olx,\oly,\olz\in\Bbbk[V]$
      分别映为$t^3,t^4,t^5\in\Bbbk[\bbA^1]$.
      \item $\fai$的逆映射只能是
      $\psi:V\to\bbA^1,\,(x,y,z)\mapsto\frac yx$,
      显然不能表示成多项式的形式.
      故$\fai$不是$\bbA^1$与$V$的同构.
    \end{enumerate}
  \end{proof}

  \begin{prob}\label{ex-15-1-27}
    设$\fai:V\to W$为仿射代数集之间的态射.
    若$W'\subseteq W$为仿射代数子集,
    证明: $V$的子集$V':=\fai^{-1}(W')$
    为仿射代数集.
    若$W'=\mcalZ(I)$,
    证明: $V'=\mcalZ(\faitil(I))$.
  \end{prob}

  \begin{proof}
    一方面,对任意$\bfv\in V'$,
    则$\fai(\bfv)\in W'$.
    对任意$f\in\faitil(I)$,
    则存在$g\in I$使得$f=\faitil(g)$.
    于是$f(\bfv)=\faitil(g)(\bfv)=g(\fai(\bfv))=0$.
    这就证明了$V'\subseteq\mcalZ(\faitil(I))$.
    另一方面,对任意$\bfv\in\mcalZ(\faitil(I))$,
    则对任意$g\in I$都有
    $g(\fai(\bfv))=\faitil(g)(\bfv)=0$,
    从而$\fai(\bfv)\in\mcalZ(I)=W'$,
    故$\bfv\in\fai^{-1}(W')=V'$.
    所以$\mcalZ(\faitil(I))\subseteq V'$.
    综上,有$V'=\mcalZ(\faitil(I))$.
    特别地,$V'$也为仿射代数集.
  \end{proof}

  \begin{prob}\label{ex-15-1-28}
    设$V,W$为仿射代数集,
    证明: $V\times W$也是仿射代数集,并且
    \[\Bbbk[V\times W]\cong\Bbbk[V]\ten_{\Bbbk}\Bbbk[W].\]
  \end{prob}

  \begin{proof}
    设$V=\mcalZ(I)\subseteq\bbA^m$, 以及
    $W=\mcalZ(J)\subseteq\bbA^n$.
    其中$I,J$分别是环$\Bbbk[x_1,x_2,...,x_m]$
    与$\Bbbk[y_1,y_2,...,y_n]$
    的理想.则显然有
    \[V\times W=\mcalZ(I\cup J)\subseteq \bbA^{m+n}\]
    其中视$I,J$视为$\Bbbk[x_1,x_2,...,x_m;\,y_1,y_2,...,y_n]$的子集.
    从而$V\times W$为仿射代数集.
    容易验证
    \begin{eqnarray*}
      \phi:\Bbbk[V]\ten_{\Bbbk}\Bbbk[W]&\to&\Bbbk[V\times W]\\
      f(\bfx)\ten g(\bfy)&\mapsto& f(\bfx)g(\bfy)
    \end{eqnarray*}
    为良定的$\Bbbk$-代数同态,并且显然为满同态.
    只需再验证$\phi$为单同态.
    设$h\in\ker\phi$,
    不妨$h=\sum_{i=1}^{s}f_i\ten g_i$,
    其中$f_i\in\Bbbk[V],\,g_i\in\Bbbk[W]$,
    并且不妨$\{g_1,g_2,...,g_s\}$在$\Bbbk[W]$当中$\Bbbk$-线性无关.
    任意取定$\bfx\in V$,
    考虑$h_{\bfx}(\bfy):=\sum_{i=1}^{s}f_i(\bfx)\ten g_i(\bfy)\in\Bbbk[W]$.
    由$\phi(h)=0\in\Bbbk[V\times W]$
    可知$h_{\bfx}=0\in\Bbbk[W]$.
    而$\{g_1,g_2,...,g_s\}$在$\Bbbk[W]$当中$\Bbbk$-线性无关,
    这迫使$f_1(\bfx)=f_2(\bfx)=\cdots=f_s(\bfx)=0$.
    这对任意$\bfx\in V$都成立.
    因此$f_1,f_2,...,f_s=0\in\Bbbk[V]$.
    从而$h=0\in\Bbbk[V]\ten_{\bfk}\Bbbk[W]$.
    从而$\phi$为单同态.
    综上,$\phi$诱导$\Bbbk$-代数同构
    \[\Bbbk[V\times W]\cong\Bbbk[V]\ten_{\Bbbk}\Bbbk[W].\]
  \end{proof}
\end{enumerate}

在以下$7$题,我们引入$R$-模$M$的相伴素理想的概念.
(也可见$\ref{section-15-4}$节的\ref{ex-15-4-30}题至
\ref{ex-15-4-40}题,以及\ref{section-15-5}节的
\ref{ex-15-5-25}题至\ref{ex-15-5-30}题).
设$M$为$R$-模,环$R$的素理想$\mfkp$称为$M$的
\newdef{相伴素理想}{associated prime},
如果存在$m\in M$使得$\mfkp=\Ann(m)$.
相伴素理想也被叫做\newdef{刺客}{assassin}.
$M$的相伴素理想之全体记作$\Ass_R(M)$.
容易知道,素理想$\mfkp\in\Ass_R(M)$当且仅当
存在$M$的某个同构于$R/\mfkp$的循环子模$Rm$.
设$I$为$R$的理想,称$\Ass_R(R/I)$中的元素为理想$I$的相伴素理想.
\remark{注意这里有术语混用,
$I$\textbf{作为$R$-模}的相伴素理想$\Ass_R(I)$
与\textbf{作为$R$的理想}的相伴素理想$\Ass_R(R/I)$一般不同.
}

\begin{enumerate}[resume]
  \begin{prob}\label{ex-15-1-29}
    设$R=\bbZ$, $M=\bbZ/n\bbZ$为$R$-模.求$\Ass_R(M)$.
  \end{prob}

  \begin{solution}
  $(p)\in\Ass_{\bbZ}(\bbZ/n\bbZ)$
  当且仅当存在$\bbZ$-模嵌入
  $\bbZ/p\bbZ\inj \bbZ/n\bbZ$,
  从而$p|n$.
  因此$\Ass_{\bbZ}(\bbZ/n\bbZ)=\Bigset{(p)}{p\text{为$n$的素因子}}$.
  \end{solution}

  \begin{prob}\label{ex-15-1-30}
    设$R$-模$M$为一族$R$-子模$M_i, (i\in\mcalA)$的并,
    证明:$\Ass_R(M)=\bigcup_{i\in\mcalA}\Ass_R(M_i)$.
  \end{prob}

  \begin{proof}
    注意$\Bigset{m\in M}{\text{$\Ann(m)$为$R$的素理想}}
    =\bigcup_{i\in\mcalA}
    \Bigset{m\in M_i}{\text{$\Ann(m)$为$R$的素理想}}$.
    从而由相伴素理想的定义易得.
  \end{proof}

  \begin{prob}\label{ex-15-1-31}
    对于$R$-模$M$,
    设素理想$\mfkp\in\Ass_R(M)$,
    $m\in M$使得$\mfkp=\Ann(m)$.
    证明:对任意$0\neq m'\in Rm\subseteq M$,
    都有$\Ann(m')=\mfkp$.
    并由此得出$\Ass_R(R/\mfkp)=\{\mfkp\}$.
    \hint{注意$R/\mfkp$是整环.}
  \end{prob}

  \begin{proof}\,
  \begin{itemize}
    \item
    设$m'=rm$,其中$0\neq r\in R$.
    由$m'\neq 0$可知$r\not\in\mfkp$.
    则显然$\mfkp=\Ann(m)\subseteq\Ann(rm)$.
    另一方面,对任意$s\in\Ann(rm)$,
    则$sr\in\Ann(m)=\mfkp$.
    注意$\mfkp$为素理想且$r\not\in\mfkp$,
    从而必有$s\in\mfkp=\Ann(m)$.
    综上则有$\Ann(m)=\Ann(m')$.

    \item 设$\mfkp$为$R$的素理想,
    视$R/\mfkp$为$R$-模.
    考虑$\overline{1}\in R/\mfkp$,
    易验证$\Ann(\overline{1})=\mfkp$.
    从而由刚才的结论,易知对任意$0\neq\olm\in R/\mfkp$
    都有$\Ann(\olm)=\mfkp$.
    故$\Ass_R(R/\mfkp)=\{\mfkp\}$.
  \end{itemize}
  \end{proof}

  \begin{prob}\label{ex-15-1-32}
    设$M$为$R$-模,
    $\mfkp$为$\mcalS:=\Bigset{\Ann(m)}{0\neq m\in M}$
    的一个(包含偏序)极大元,证明: $\mfkp$为$R$的素理想.
    \hint{若$\mfkp=\Ann(m)$且$ab\in\mfkp$,
    证明$bm\neq 0$蕴含$\Ann(m)\subseteq\Ann(bm)$.
    之后利用$\mfkp=\Ann(m)$的极大性得出$a\in\Ann(bm)=\mfkp$.}
  \end{prob}

  \begin{proof}
    设$\mfkp=\Ann(m)$为$\mcalS$的一个极大元.
    设$a,b\in R$使得$ab\in\mfkp$,
    则$abm=0$,即$a\in\Ann(bm)$.
    如果$b\not\in\mfkp$,
    则$bm\neq 0$,
    从而$\Ann(bm)\in\mcalS$.
    而显然$\mfkp=\Ann(m)\subseteq\Ann(bm)$,
    于是$\mfkp$的极大性迫使$\mfkp=\Ann(bm)$.
    故$a\in\Ann(bm)=\mfkp$.
    这就证明了$\mfkp$为素理想.
  \end{proof}

  \begin{prob}\label{ex-15-1-33}
    设$R$为Noether环, $M\neq 0$为$R$-模.
    证明: $\Ass_R(M)\neq\vkong$.
    \hint{利用上题.}
  \end{prob}

  \begin{proof}
    沿用上题记号,
    注意$\mcalS:=\Bigset{\Ann(m)}{0\neq m\in M}$
    为$R$的一族理想,且$\mcalS\neq\vkong$.
    $R$为Noether环保证了$\mcalS$之中必有极大元
    \remark{这利用了Noether环的等价定义,并不依赖Zorn引理},
    由上题知该极大元必为素理想,
    从而属于$\Ass_R(M)$.
    因此$\Ass_R(M)$非空.
  \end{proof}

  \begin{prob}\label{ex-15-1-34}
    设$R$-模正合列$0\to L\to M\to N\to 0$,
    证明包含关系
    \footnote{原书将(*)式最左边的$\Ass_R(L)\subseteq\Ass_R(M)$
    误写为$\Ass_R(N)\subseteq\Ass_R(M)$.}
    \[
      \Ass_R(L)\subseteq\Ass_R(M)
      \subseteq\Ass_R(N)\cup\Ass_R(L).\eqno{(*)}
    \]
    并举例说明上述两个包含都可以是真包含.
    \hint{若$Rm\cong R/\mfkp$,
    证明$Rm\cap L=0$蕴含$p\in\Ass_R(N)$;
    而当$Rm\cap L\neq 0$时必有$\mfkp\in\Ass_R(L)$,
    利用习题\ref{ex-15-1-31}.
    至于真包含关系,考虑$n\bbZ\subseteq\bbZ$.}
  \end{prob}

  \begin{proof}\,
  \begin{itemize}
    \item
    对于$M$的子模$L$,
    设$\mfkp\in\Ass_R(L)$.
    则存在$l\in L$使得$\mfkp=\Ann(l)$.
    只需注意$l\in M$也成立,从而$\mfkp\in\Ass_R(M)$.
    这就证明了$\Ass_R(L)\subseteq\Ass_R(M)$.
    又设$\mfkq\in\Ass_R(L)$,
    取$m\in M$使得$\mfkq=\Ann(m)$.
    如果$Rm\cap L\neq 0$,
    则任取$Rm\cap L$的非零元$m'$,
    则由习题$\ref{ex-15-1-31}$可知$\Ann(m')=\mfkq$,
    从而$q\in\Ass_R(L)$.
    而如果$Rm\cap L=0$,
    则考虑$m$在$N\cong M/L$当中的像$\olm\neq 0$.
    由$Rm\cap L=0$易验证$\Ann(\olm)=q$,
    这表明$\mfkq\in\Ass_R(N)$.
    综上所述,成立
    $\Ass_R(M)
      \subseteq\Ass_R(N)\cup\Ass_R(L)$.
    \item 设$p,q$为不同素数,
    考虑$\bbZ$-模$\bbZ/p\bbZ\inj\bbZ/pq\bbZ$,
    而由习题\ref{ex-9-6-29}可知$\Ass_{\bbZ}(\bbZ/p\bbZ)
    \subsetneqq\Ass_{\bbZ}(\bbZ/pq\bbZ)$,
    这是一个真包含关系.
    再考虑正合列
    \[0\to p\bbZ\to\bbZ\to\bbZ/p\bbZ \]
    易验证$\{(0)\}=\Ass_{\bbZ}\bbZ\subsetneqq
    \Ass_{\bbZ}(p\bbZ)\cup\Ass_{\bbZ}(\bbZ/p\bbZ)
    =\{(0),(p)\}$.这是另一个真包含关系.
  \end{itemize}
  \end{proof}

  \begin{prob}\label{ex-15-1-35}
    设$M$为$R$-模, $\mcalS$为$\Ass_R(M)$的一个子集.
    证明:存在$M$的子模$N$使得$\Ass_R(N)=\mcalS$,
    并且$\Ass_R(M/N)=\Ass_R(M)\setminus\mcalS$.
    \hint{考虑使得$\Ass_R(N')\subseteq\mcalS$的$M$的子模$N'$构成的集合,
    利用习题\ref{ex-15-1-30}以及Zorn引理断言存在该集合中的极大元$N$.
    若$\mfkp\in\Ass_R(M/N)$,
    则存在子模$M'/N\cong R/\mfkp$.
    利用之前的习题去证明
    $\Ass_R(M')\subseteq\Ass_R(N)\cup\{\mfkp\}$,
    再利用$N$的极大性去证明$\mfkp\in\Ass_R(M)\setminus\mcalS$.
    从而$\Ass_R(M/N)\subseteq\Ass_R(M)\setminus\mcalS$,
    并且$\Ass_R(N)\subseteq\mcalS$.
    再次利用之前习题即可.
    }
  \end{prob}

  \begin{proof}


    考虑集合$\mcalL:=\Bigset{N\text{为$M$的子模}}{\Ass_R(N)\subseteq\mcalS}$.
    则对$\mcalL$中的任一(包含偏序)链$\{N_i\}$,
    由习题\ref{ex-15-1-30}可知
    $\Ass_R(\bigcup_i N_i)=\bigcup_i\Ass_R(N_i)\subseteq\mcalS$,
    从而$\bigcup_i N_i\in\mcalL$.
    从而由Zorn引理知$\mcalL$中存在(包含偏序)极大元.
    设$N$为$\mcalL$的一个极大元.则$\Ass_R(N)\subseteq\mcalS$.
    而对于任意$\mfkp\in\Ass_R(M/N)$,
    则存在$M$的子模$M'\supseteq N$使得$R/\mfkp\cong M'/N$.
    注意$R$-模正合列
    \[0\to N\to M'\to M'/N\to 0,\]
    从而由习题\ref{ex-15-1-34}可知
    $\Ass_R(M')\subseteq\Ass_R(N)\cup\Ass_R(M'/N)=\Ass_R(N)\cup\{\mfkp\}$.
    又因为$M'\supseteq N$,
    从而$N$在$\mcalL$中的极大性迫使$\Ass_R(M')\not\subseteq\mcalS$,
    从而只能$\mfkp\in\Ass_R(M')$并且$\mfkp\not\in\mcalS$.
    而$\Ass_R(M')\subseteq\Ass_R(M)$,
    这就证明了$\mfkp\in\Ass_R(M)\setminus\mcalS$.
    因此$\Ass_R(M/N)\subseteq\Ass_R(M)\setminus\mcalS$.
    而由$0\to N\to M\to M/N\to 0$知
    \[\Ass_R(M)\subseteq\Ass_R(N)\cup\Ass_R(M/N)
    \subseteq \mcalS\cup(\Ass_R(M)\setminus\mcalS)=\Ass_R(M)\]
    这迫使$\Ass_R(N)=\mcalS$并且$\Ass_R(M/N)=\Ass_R(M)\setminus\mcalS$.
  \end{proof}

\end{enumerate}

设$M$为有限生成$R$-模,
$\{m_1,m_2,...,m_n\}$为$M$的一组$R$-模生成元.
记$\mcalF_R(M)$为$R$的由以下元素生成的理想:
\[
  \Bigset{
    \det A}
  {A=(r_{ij})_{n\times n}\in R^{n\times n},
  \sum_{j=1}^{n}r_{ij}m_j=0,\,\forall 1\leq i\leq n.}
\]
称$\mcalF_R(M)$为$M$的\newdef{Fitting理想}{Fitting ideal},
也俗称"行列式理想".注意到上述矩阵$A$的每一行都
可视为为$m_1,m_2,...,m_n$的一个生成关系,
故如此的矩阵$A$也称为$M$的"关系矩阵".
以下$5$题介绍Fitting理想的基本性质.
\remark{上述方式定义的$\mcalF_R(M)$
似乎与生成元组$\{m_1,m_2,...,m_n\}$的选取有关.
后面的习题将说明此定义的良定性.}

\begin{enumerate}[resume]
  \begin{prob}\label{ex-15-1-36}
  设$M$为有限生成$R$-模,
  $\{m_1,m_2,...,m_n\}$为$M$的一组$R$-模生成元.
    \begin{enumerate}
      \item 证明$\mcalF_R(M)$
      也可由所有$A\in R^{p\times r}$
      的所有的$n$阶子式生成,
      其中$p\geq n$,
      并且$B$的每一行都是生成元$m_1,m_2,...,m_n$的生成关系.

      \item 设$A$为给定的$p\times n$的$R$-矩阵, $p\geq n$.
      $A'$为$A$经过某些初等行,列变换所得到的矩阵.
      证明: $R$的由$A$的所有$n$阶子式生成的理想
      与由$A'$的所有$n$阶子式生成的$R$理想相等.
    \end{enumerate}
  \end{prob}

  \begin{proof}
    (a)是显然的.至于(b),
    对于$p\times n$的$R$-矩阵$A$,
    记$S_A$为$A$的所有$n$阶子式构成的集合.
    记$S_A$生成的$R$的理想为$I_A$.
    注意对换两行,
    把其中一行倍乘$R$中的可逆元
    (以及关于列的类似操作)
    在相差$R$中某可逆元倍的意义下不改变集合$S_A$,
    从而不改变$S_A$生成的$R$的理想$I_A$.
    至于将一行倍加到另一行的情形,
    这对应着将$S_A$中的某元素倍加到其余某些元素上,
    这也不改变$S_A$生成的理想$I_A$.
    从而得证.
  \end{proof}

  \begin{prob}\label{ex-15-1-37}
    设$\{m_1,m_2,...,m_n\}$与$\{m'_1,m'_2,...,m'_{n'}\}$
    为$R$-模$M$的两组生成元.
    设$\mcalF$为$M$关于$\{m_1,m_2,...,m_n\}$的Fitting理想,
    $\mcalF'$为$M$关于$\{m_1,m_2,...,m_n;
    m'_1,m'_2,...,m'_{n'}\}$的Fitting理想.
    我们将证明$\mcalF=\mcalF'$,
    进而断言$\mcalF_R(M)$与$M$的生成元组选取无关.
    \begin{enumerate}
      \item 证明:对任意$1\leq s\leq n'$,
      存在$a_{s'1},a_{s'2},...,a_{s'n}\in R$,
      使得$m_s'=\sum_{i=1}^{n}a_{s'i}m_i$.
      从而断言: $(-a_{s'1},-a_{s'2},...,-a_{s'n};0,...,0,1,0,...,0)$
      为$m_1,m_2,...,m_n; m'_1,m'_2,...,m'_{n'}$的一个生成关系.

      \item 设$A\in R^{n\times n}$,并且$A$的每一行都是$m_1,m_2,...,m_n$的生成关系;
      再设$A'\in R^{(n+n')\times(n+n')}$,
      并且$A'$的每一行都是$m_1,m_2,...,m_n; m'_1,m'_2,...,m'_{n'}$的生成关系.
      证明: $\det A=\det A'$.
      从而推出$\mcalF\subseteq\mcalF'$.
      \hint{利用上一小问,对$A'$适当初等变换,使得变换之后为分块上三角阵,
      并且$n\times n$左上分块为$A$,右下分块为$n'\times n'$单位阵.}
      \footnote{译者注:此小问有误,貌似缺少条件.
      但不影响后面.}

      \item 证明: $\mcalF'\subseteq\mcalF$,
      从而推出$\mcalF=\mcalF'$.
      \hint{利用之前的习题.}

      \item 证明: Fitting理想$\mcalF_R(M)$与$M$的生成元组选取无关.
    \end{enumerate}
  \end{prob}

  \begin{proof}\,
    \begin{enumerate}
      \item 这是因为$\{m_1,m_2,...,m_n\}$是$M$的一组
      $R$-模生成元,且$m_s'\in M$,
      从而$m_s'$能够表示为$\{m_1,m_2,...,m_n\}$的$R$-线性组合.
      之后显然.
      \item \remark{这一小问显然不对.设$A,A'$满足题设,且$\det A=\det A'$,
      则任取$1\neq r\in R$,将$A$的某一行倍乘$r$,
      得到的矩阵$B$的各行依然是$\{m_1,m_2,...,m_n\}$的生成关系.
      从而满足题设,故应该有$\det B=\det A'$.
      但$\det B=r\det A=r\det A'$.如果$R$是整环且$\det A\neq 0$,
      这是不可能的.}
      \item 无视上一问,直接证明如下.
      一方面,对任意$\det A\in\mcalF$,
      其中$A$为$n\times n$阶$R$-方阵,
      且各行都是$\{m_1,m_2,...,m_n\}$的生成关系,
      则考虑$A':=
      \begin{pmatrix}
        A & O \\
        O & O
      \end{pmatrix}$
      为$(n+n')\times (n+n')$阶$R$-方阵,
      且各行显然都是$\{m_1,...,m_n; m_1',...,m_{n'}'\}$
      的生成关系,
      故$\det A=\det A'\in\mcalF'$,
      所以$\mcalF\subseteq\mcalF'$.
      另一方面,
      对任意$\det A'\in\mcalF'$,
      其中$A'=
      \begin{pmatrix}
        A'_{11} & A'_{12} \\
        A'_{21} & A'_{22}
      \end{pmatrix}$
      为$(n+n')\times(n+n')$阶$R$-矩阵,
      并且每一行都为$\{m_1,...,m_n; m_1',...,m_{n'}'\}$
      的生成关系.
      则考虑矩阵$B=
      \begin{pmatrix}
        A'_{11} & A'_{12} \\
        A'_{21} & A'_{22} \\
        K & I
      \end{pmatrix}$,
      其中$K:=-(a_{ij})_{1\leq i\leq n'\atop 1\leq j\leq n
      }$,其中$a_{ij}$的含义见本题(a)问.
      则$\det A'\in I_B$,
      其中$I_B$为$B$的所有的$(n+n')$阶子式生成的理想.
      对矩阵$B$适当做初等行变换,
      可将$B$化为
      $B'=
      \begin{pmatrix}
        B_1 & O \\
        B_2 & O \\
        K & I
      \end{pmatrix}$,
      易知$B_1,B_2$的每一行的前$n$列都是
      $\{m_1,m_2,...,m_n\}$的生成关系.
      由上一题知$I_B=I_{B'}$.
      而易知$B'$的$(n+n')$阶子式必为
      $\begin{pmatrix}
         B_1 \\
         B_2
       \end{pmatrix}$
      的某个$n$阶子式,因此综上有
      \[\det A'\in I_B=I_{B'}\subseteq
      I_{{B_1\choose B_2}}\subseteq \mcalF\]
      这就证明了$\mcalF'\subseteq\mcalF$.
      综上所述, $\mcalF'=\mcalF$.

      \item 若再记关于生成元$\{m_1',m_2',...,m_{n'}'\}$
      的Fitting理想为$\mcalF''$,
      则同理也有$\mcalF''=\mcalF'$.
      因此$\mcalF=\mcalF''$.
      这表明Fitting理想与$M$的生成元选取无关.
    \end{enumerate}

  \end{proof}

  \begin{prob}\label{ex-15-1-38}
    本题中所有的$R$-模都默认是有限生成的.
    \begin{enumerate}
      \item 若$R$-模$M$可被$n$个元素生成,证明:
      $\Ann(M)^n\subseteq\mcalF_R(M)\subseteq\Ann(M)$,
      其中$\Ann(M)$为$R$-模$M$的湮灭子.
      \hint{若$A$为$M$的$n\times n$关系矩阵,
      则$AX=0$,其中$X$为$M$的那$n$个生成元排成的列向量.
      将$AX=0$左乘$A$的伴随矩阵,断言$\det A\in\Ann(M)$.}
      \item 设$M_1,M_2$为$R$-模,证明:
      $\mcalF_R(M_1\oplus M_2)=\mcalF_R(M_1)\mcalF_R(M_2)$.
      \item 设$R$-模$M=(R/I_1)\oplus(R/I_2)
      \oplus\cdots\oplus(R/I_n)$
      为$n$个循环$R$-模的直和,
      其中$I_1,...,I_n$为$R$的理想.
      证明: $\mcalF_R(M)=I_1I_2\cdots I_n$.
      \item 若$R=\bbZ$, $M$为有限生成Abel群,
      证明: 当$M$为无限群时,
      $\mcalF_{\bbZ}(M)=0$;
      当$M$为有限群时, $\mcalF_{\bbZ}(M)=|M|\bbZ$.
      \item 设$I$为$R$的理想,$M$为$R$-模.
      证明: $\mcalF_R(M)$在$R/I$当中的像等于$\mcalF_{R/I}(M/IM)$.
      \item 记号同上,证明:
      $\mcalF_R(M/IM)\subseteq\mcalF_R(M)+I\subseteq R$.
      \item 若$\fai:M\to M'$为$R$-模满同态.
      证明: $\mcalF_R(M)\subseteq\mcalF_R(M')$.
      \item 若$0\to L\to M\to N\to 0$为$R$-模短正合列,
      证明: $\mcalF_R(L)\mcalF_R(N)\subseteq\mcalF_R(M)$.
      \item 设$R=\Bbbk[x,y,z]$,其中$\Bbbk$为域.
      设$M=R/(x,y^2,yz,z^2)$,
      $L=(x,y,z)/(x,y^2,yz,z^2)$为$M$的$R$-子模.
      证明: $\mcalF_R(M)=(x,y^2,yz,z^2)$,
      以及$\mcalF_R(L)=(x,y,z)^2$.
      \remark{一般地,子模$L\subseteq M$的Fitting理想未必包含$M$的Fitting理想.}
    \end{enumerate}
  \end{prob}

  \begin{proof}\,
    \begin{enumerate}
      \item 设$\{m_1,m_2,...,m_n\}$为$M$的一组生成元.
      对任意$a=a_1a_2\cdots a_n\in\Ann(M)^n$,
      其中$a_k\in\Ann(M)$, $1\leq k\leq n$,
      则矩阵$\Lmd:=\begin{pmatrix}
            a_1 &  &  &  \\
             & a_2 &  &  \\
             &  & \ddots &  \\
             &  &  & a_n
          \end{pmatrix}$
      是关于$\{m_1,m_2,...,m_n\}$的生成关系矩阵,
      因此$a=\det \Lmd\in\mcalF_R(M)$.
      这就证明了$\Ann(M)^n\subseteq\mcalF_R(M)$.
      另一方面,对任意$\det A\in\mcalF_R(M)$,
      其中方阵$A$为关于$\{m_1,m_2,...,m_n\}$
      的生成关系矩阵.
      考虑$A$的伴随方阵$A^*$,
      即$A^*$的矩阵元为$A$的相应的代数余子式,
      则$A^*A=\det A\cdot I_{n}$也为生成关系矩阵,
      因此必有$\det A\in\Ann(M)$.
      这就证明了$\mcalF_R(M)\subseteq\Ann(M)$.

      \item 设$G:=\{m_1,m_2,...,m_n\}$
      与$G':=\{m_{1}',m_2',...,m_{n'}'\}$
      分别为$M_1,M_2$的一组$R$-模生成元.
      则$G'':=\{m_1,m_2,...,m_n\,;m_{1}',m_2',...,m_{n'}'\}$
      为$M\oplus M'$的一组$R$-模生成元.
      一方面,对任意的关于$G$的生成关系方阵$A$以及关于$G'$
      的生成关系方阵$A'$,
      则分块矩阵$ A'':=
      \begin{pmatrix}
        A &  \\
          & A'
      \end{pmatrix}$
      为关于$G''$的生成关系方阵.
      所以$\det A\cdot\det A'=\det A''\in\mcalF_R(M_1\oplus M_2)$.
      这就证明了$\mcalF_R(M_1)\mcalF_R(M_2)
      \subseteq\mcalF_R(M_1\oplus M_2)$.
      另一方面,对任意的关于$G''$的生成关系方阵$B$,
      则$B$的任意一行的前$n$列,后$n'$列分别为$G,G'$的生成关系.
      于是将行列式$\det B$按前$n$列作Laplace展开,
      展开式中的每一项都位于$\mcalF_R(M_1)\mcalF_R(M_2)$.
      这就证明了$\mcalF_R(M_1\oplus M_2)
      \subseteq\mcalF_R(M_1)\mcalF_R(M_2)$.
      综上所述,有$\mcalF_R(M_1\oplus M_2)
      =\mcalF_R(M_1)\mcalF_R(M_2)$.

      \item 对于环$R$的理想$I$,
      视商环$R/I$为$R$-模,
      则$R/I$可由一个元素$\overline{1}$生成,
      其中$1\in R$为乘法单位元.
      设$1\times 1$阶$R$-方阵(其实就是$R$中的元素)$r$
      为$R/I$的关于$\{\overline{1}\}$的生成关系矩阵,
      则显然$r\in I$.
      这表明$\mcalF_R(R/I)=I$.
      再利用(b)的结论即可.

      \item 注意有限生成Abel群的结构.
      如果$M$为无限群,则$\bbZ$必为$M$的一个直和项.
      而由(c)知$\mcalF_{\bbZ}(\bbZ)=0$,
      因此再由(b)易知$\mcalF_{\bbZ}(M)=0$.
      而当$M$为有限群时,$M$必形如
      $\bigoplus_{k=1}^n\bbZ/a_k\bbZ$,
      其中$a_k$为正整数.
      而$\mcalF_{\bbZ}(\bbZ/a_k\bbZ)=a_k\bbZ$,从而
      $
        \mcalF_{\bbZ}(M)
      =\prod_{k=1}^{n}
        (a_k\bbZ)
      =\left(
        \prod_{k=1}^{n}
          a_k
       \right)\bbZ
      =|M|\bbZ.
      $
      \item 设$G:=\{m_1,m_2,...,m_n\}$为$M$的一组$R$-模生成元.
      则易验证$\overline{G}:=\{\olm_1,\olm_2,...,\olm_n\}$
      为$M/IM$的$R/I$-模生成元,
      其中$\olm_i$为$m_i$在$M/IM$中的像.
      一方面,设$n$阶$R$-方阵$A=(a_{ij})$为关于$G$的生成关系矩阵,
      则$R/I$-方阵$\olA=(\ola_{ij})$
      为关于$\olG$的生成关系矩阵,
      其中矩阵元$\ola_{ij}$为$a_{ij}$在$R/I$中的像.
      从而$\det\olA\in\mcalF_{R/I}(M/IM)\subseteq R/I$.
      这就证明了$\overline{\mcalF_R(M)}\subseteq\mcalF_{R/I}(M/IM)$.

      另一方面,对任意$\olB=(\olb_{ij})$
      为关于$G'$的生成关系方阵,其矩阵元$\olb_{ij}\in R/I$.
      任取$b_{ij}\in R$为$\olb_{ij}$的一个代表元.
      则对于任意$1\leq i\leq n$,
      有$\sum_{j=1}^{n}b_{ij}m_j\in IM$.
      注意$IM$当中的元素必可以写成$\{m_1,m_2,...,m_n\}$的$I$-线性组合,
      因此存在$\lmd_{ij}\in I$使得
      \[\sum_{j=1}^{n}b_{ij}m_j=\sum_{j=1}^{n}\lmd_{ij}m_j\]
      这对任意$1\leq i,j\leq n$都成立.
      记方阵$L=(\lmd_{ij})$,
      则$B-L$为关于$G$的生成关系矩阵,
      并且易知$\det B\equiv\det (B-L)\mod I$.
      这就证明了$\mcalF_{R/I}(M/IM)\subseteq\overline{\mcalF_R(M)}$.
      综上证毕.

      \item 记号同上一小问,
      设$R$-方阵$B=(b_{ij})$为$G'$的一个生成关系矩阵,
      则$\sum_{j=1}^{n}b_{ij}m_j\in IM$.
      仿照上一问也有类似的$L$,使得$L$的各矩阵元都属于$I$,
      并且$B-L$为$G$的生成关系.
      再注意到$\det B\equiv\det(B-L)\mod I$,
      从而易知$\mcalF_R(M/IM)\subseteq\mcalF_R(M)+I$.证毕.

      \item 设$G:=\{m_1,m_2,...,m_n\}$
      为$M$的一组$R$-生成元,
      由$\fai:M\to M'$为满同态可知
      $G':=\{\fai(m_1),\fai(m_2),...,\fai(m_n)\}$
      为$M'$的$R$-生成元.
      则对于$G$的生成关系矩阵$A$,
      易验证$A$也是$G'$的生成关系,所以$\det A\in\mcalF_R(M')$.
      因此$\mcalF_R(M)\subseteq\mcalF_R(M')$.

      \item 设$G:=\{g_1,g_2,...,g_m\}$为$L$
      的一组$R$-模生成元,
      $H:=\{\olh_1,\olh_2,...,\olh_n\}$
      为$N\cong M/L$的一组$R$-生成元.
      任取$\olh_j$在$M$中的代表元$h_j$,
      则易知$K:=\{g_1,...,g_m\,;h_1,h_2,...,h_n\}$
      为$M$的$R$-生成元.
      现在,设$A,B$分别为$G,H$的生成关系方阵,
      易验证存在分块方阵$C=
      \begin{pmatrix}
        A &   \\
        * & B
      \end{pmatrix}$
      为关于$K$的生成关系方阵.
      从而$\det A\det B=\det C\in\mcalF_R(M)$.
      这就证明了$\mcalF_R(L)\mcalF_R(N)\subseteq\mcalF_R(M)$.

      \item 注意$M$为循环$R$-模,
      从而由(c)知显然有$\mcalF_R(M)=(x,y^2,yz,z^2)$.
      至于$L=(x,y,z)/(x,y^2,yz,z^2)$,
      则显然$G:=\{\oly,\olz\}$为$L$的一组$R$-模生成元.
      设$r_1,r_2\in R$满足$r_1\oly+r_2\olz=\overline{0}\in L$,
      则易知必有$r_1,r_2\in(x,y,z)$,
      从而用定义直接计算得$\mcalF_R(L)=(x,y,z)^2$.

    \end{enumerate}
  \end{proof}

  \begin{prob}\label{ex-15-1-39}
    设$\fai:R^n\to M$为$R$-模满同态
    (特别地,$M$可由$n$个元素有限生成).
    记$L=\ker\fai$.
    证明: $\mcalF_R(M)=
    \im\left(\bigwedge\nolimits^n(L)
    \to\bigwedge\nolimits^n(R^n)\right)$.
    其中$\faitil:
    \bigwedge\nolimits^n(L)\to\bigwedge\nolimits^n(R^n)\cong R$
    由包含映射$L\inj R^n$所诱导.
  \end{prob}

  \begin{proof}
    设$\{\bfe_1,\bfe_2,...,\bfe_n\}$为$R^n$的标准基,
    则$G:=\{\fai(\bfe_1),\fai(\bfe_2),
    ...,\fai(\bfe_n)\}$
    为$M$的一组$R$-生成元.
    对于行向量$\bfv_i\in R^n$, $1\leq i\leq n$,
    若$\bfv_i$为关于$G$的生成关系,
    则$\bfv_i\in\ker\fai=L$.
    再由行列式与外代数的关系:
    \[\det(\bfv_1,\bfv_2,...,\bfv_n)^T
    =\faitil(\bfv_1\wedge\bfv_2\wedge\cdots\wedge\bfv_n)\]
    立刻得证.
  \end{proof}

  \begin{prob}\label{ex-15-1-40}
    设$\fai:R\to S$为环同态.
    $M$为有限生成$R$-模,
    证明: $\mcalF_S(M\ten_RS)\cong\mcalF_R(M)\ten_RS$.
  \end{prob}
\end{enumerate}

\begin{proof}
  取$R$-模满同态$f: R^n\to M$,记$L=\ker f\subseteq R^n$,
  则有$R$-模右短正合列
  \[
    L\to R^n\to M\to 0
  \]
  注意$-\ten_RS$的右正合性,从而有$R$-模右短正合列
  \[
    \xymatrix{
       L\ten_RS
         \ar[r]
      &R^n\ten_RS
         \ar[r]
         \ar@{=}[d]^{\cong}
      &M\ten_RS
         \ar[r]
      &0
    \\
      &S^n
      &&
    }
  \]
  显然这也是作为$S$-模的右正合列.从而由上题可知
  \begin{eqnarray*}
    \mcalF_S(M\ten_RS)
  &=&
    \im\left(
      \bigwedge\nolimits^n(L\ten_RS)\to S
    \right)\\
  &=&
    \im\left(
      \bigwedge\nolimits^n(L)\ten_RS\to S
    \right)\\
  &\cong&
  \im\left(
      \bigwedge\nolimits^n(L)\to R
    \right)\ten_RS\\
  &=&
    \mcalF_R(M)\ten_RS.
  \end{eqnarray*}
\end{proof}

\end{document} 